\documentclass[10pt,a4paper]{article}

% ----------------- Paquetes -------------------%
\usepackage[T1]{fontenc}
\usepackage[utf8]{inputenc}
\usepackage[a4paper,margin=2.5cm]{geometry}
\usepackage{mathptmx}             
\usepackage{changepage} 
\usepackage{setspace}
\usepackage[spanish]{babel}
\usepackage[labelfont=bf,labelsep=space]{caption}
\usepackage{amssymb}
\usepackage{amsmath}
\usepackage{ebgaramond}
\usepackage{placeins}
\usepackage{etoolbox}
\usepackage{setspace}
\usepackage{parskip} 
\usepackage{tcolorbox}
\usepackage{needspace}
\usepackage{xparse}
\usepackage{ocgx2}
\usepackage{hyperref}
\usepackage{hypcap}
\usepackage{soul}\sethlcolor{yellow}% resaltado amarillo: \hl{texto} (Panel TCEE, ⌃H)


%%%%%%%%%%%%%%%%%%%%%%%%%%%%%%%%%%%%%%%%%%%%%%%%%%%%%%%%%%%%%%%%
% --- Fija primero tus valores globales "buenos" ---
\setstretch{1.2}                 % Interlineado global
\setlength{\parskip}{0.7\baselineskip}
\setlength{\parindent}{0pt}
\newlength{\GlobalParskip}
\setlength{\GlobalParskip}{\parskip}

\newcounter{eqnum}[subsection]
\renewcommand{\theeqnum}{\arabic{eqnum}}
\newcounter{alphaitem}
\newcounter{numitem}

%% ------------------- Recuadros----------------------%%
\newcommand{\puntosclave}[1]{%
  \begin{tcolorbox}[
    colframe=black,
    title={Puntos clave},
    before upper={\setlength{\parindent}{0.5cm}\setlength{\parskip}{0.5\baselineskip}}
  ]
  #1
  \end{tcolorbox}
}

\newcommand{\modificaciones}[1]{
  \begin{tcolorbox}[
    colback=white,
    colframe=red,
    title={Modificaciones a realizar},
    before upper={\setlength{\parindent}{0.5cm}\setlength{\parskip}{0.5\baselineskip}}
  ]
  #1
  \end{tcolorbox}
}

%% ------------------- Listas----------------------%%
    % --- Lista alfabética ---
    \newenvironment{la}
      {\begin{list}{\alph{alphaitem})}{%
          \usecounter{alphaitem}%
          \setlength{\leftmargin}{1cm}%
          \setlength{\parskip}{3pt}% 
          \setlength{\itemsep}{0pt}% ← espacio entre ítems
          \setlength{\parsep}{0pt}%  ← párrafos dentro de un ítem
      }}
      {\end{list}}
      
    % --- Lista numérica ---
    \newenvironment{lnum}
      {\begin{list}{\arabic{numitem}.}{%
          \usecounter{numitem}%
          \setlength{\leftmargin}{1cm}%
          \setlength{\parskip}{3pt}% 
          \setlength{\itemsep}{0pt}% ← espacio entre ítems
          \setlength{\parsep}{0pt}%  ← párrafos dentro de un ítem
      }}
      {\end{list}}
      
% ------------------- Imágenes----------------------%
    \graphicspath{{Img/}}% Rutas por defecto 
    % --- Imagen adaptativa ---
    % Registros de dimensión definidos una sola vez
    \newdimen\imgcap
    \newdimen\imgmin
    \newdimen\imgmax
    \newdimen\imgremain
    \newdimen\imgh
    \newdimen\imgneed
    
    \newcommand{\imagenfit}[3]{%
      \addvspace{10pt}% separación antes
      \par\begingroup
        % Parámetros de composición (asignaciones, no \newdimen)
        \imgcap=1\baselineskip            % alto estimado de título+fuente
        \imgmin=0.15\textheight
        \imgmax=0.15\textheight
        \imgremain=\pagegoal
        \ifdim\pagegoal=\maxdimen \imgremain=0pt\fi
        \advance\imgremain by -\pagetotal
        \advance\imgremain by -\imgcap
        % Decisión de altura destino
        \ifdim\imgremain<\imgmin
          \newpage
          \imgh=0.20\textheight
        \else
          \imgh=\imgremain
          \ifdim\imgh>\imgmax \imgh=\imgmax \fi
        \fi
        % Reservar espacio para evitar cortes del bloque completo
        \imgneed=\imgh \advance\imgneed by \imgcap
        \Needspace{\imgneed}
        % Cuerpo
        \noindent\begin{minipage}{\textwidth}
          \centering
          \captionof{figure}{\textemdash\ #2}\par\vspace{0.3em}% título arriba
          \includegraphics[width=\linewidth,height=\imgh,keepaspectratio]{\detokenize{#1}}\par
          \vspace{0.3em}\small\textit{Fuente: }#3% fuente abajo
        \end{minipage}%
      \endgroup
      \par\addvspace{10pt}%
    }

%------------ Bloque de RECUADRO ESPECIAL -------------%
    \newenvironment{specialblock}[1]{%
      \begin{quote} % crea sangría a izquierda y derecha
      \small % texto más pequeño
      \noindent\textbf{#1}\\[-0.3em] % título en negrita
      \rule{\linewidth}{0.4pt}\\[0.6em]}{
      \end{quote}}
      

%-------------------------- Author Font -----------------------%
    \newcommand{\authorfont}[1]{{\fontfamily{EBGaramond-TLF}\selectfont #1}}

%-------------------------- Anexos -----------------------%
    \makeatletter
    \newcounter{anexo}
    \renewcommand{\theanexo}{\Roman{anexo}}
    \newcommand{\anexo}[1]{%
      \clearpage
      \phantomsection
      \refstepcounter{anexo}%
      \noindent\textbf{Anexo \theanexo.- }#1\par\medskip
      \addcontentsline{stc}{section}{Anexo \theanexo.- #1}%
    }
    \makeatother

    %-------------------------- Lista de figuras (personal) -----------------------%
\makeatletter
\newcommand{\MyListOfFigures}{%
  \clearpage
  \phantomsection
  {\centering\bfseries\Large Índice de figuras\par}%
  \medskip
  % Entrada en nuestro índice de contenido (stc)
  \addcontentsline{stc}{section}{Índice de figuras}%
  % Recuperación del listado (sin cabecera propia)
  \begingroup
    \parindent=0pt\parskip=2pt
    \@starttoc{lof}%
  \endgroup
}
\makeatother

%-------------------------- Bibliografía -----------------------%
\makeatletter
% Contador para las referencias
\newcounter{myref}
% Cabecera de la bibliografía, entra en el índice 'stc'
\newcommand{\MyBibliography}{%
  \clearpage
  \phantomsection
  {\centering\bfseries\Large Bibliografía y referencias\par}%
  \medskip
  \addcontentsline{stc}{section}{Bibliografía y referencias}%
  \setcounter{myref}{0}%
}
\newcommand{\MyBibItem}[4]{%
  \refstepcounter{myref}%
  \noindent[\themyref]\ #1.\ \emph{#2}. #3. #4\par
}
\makeatother

%--------- Establecer cuatro niveles de índice --------%
    \setcounter{tocdepth}{4}   
    \setcounter{secnumdepth}{4}% numera hasta \paragraph

%%%%%%%%%%%%%%%%%%%%%%%%%%%%%%%%

\makeatletter

    %--------- Interlineado Global --------%
    \edef\GlobalBaseLineStretch{\baselinestretch}
    
    %--------- Espaciado de seccinones --------%
    \renewcommand\section{\@startsection{section}
      {1}{\z@}
      {2.5ex \@plus 1ex \@minus .2ex} % espacio antes
      {1.5ex \@plus .2ex}             % espacio después
      {\normalfont\Large\bfseries}    % formato del título
    }
    \renewcommand\subsection{\@startsection{subsection}{2}{\z@}
      {3ex \@plus 0.8ex \@minus .2ex}
      {1.5ex \@plus .2ex}
      {\normalfont\large\bfseries}
    }
    \renewcommand\subsubsection{\@startsection{subsubsection}{3}{\z@}
      {3ex \@plus 0.5ex \@minus .2ex}
      {1ex \@plus .2ex}
      {\normalfont\normalsize\bfseries}
    }
    \renewcommand\paragraph{\@startsection{paragraph}{4}{\z@}%
    {-2ex \@plus -0.5ex \@minus -.2ex}% espacio antes
    {0.8ex \@plus .2ex}% espacio después
    {\normalfont\normalsize\itshape}}% ← cursiva en lugar de negrita

    %--------- Ecuaciones --------%   
    % Variables globales para almacenar títulos y notas
    \gdef\leq@current@section{}%
    \gdef\leq@current@subsection{}%
    \gdef\leq@last@printed@section{}%
    \gdef\leq@last@printed@subsection{}%
    
    % Comando para añadir nota (invisible en el cuerpo)
    \newcommand{\eqnote}[1]{%
      \addcontentsline{leq}{leqnote}{#1}%
    }
    
    % Comando para ecuaciones
    \newcommand{\eqblock}[2]{%
      % Verificar si necesitamos imprimir el título de sección
      \ifx\leq@current@section\leq@last@printed@section\else
        \addcontentsline{leq}{leqsection}{\leq@current@section}%
        \global\let\leq@last@printed@section\leq@current@section
        \gdef\leq@last@printed@subsection{}% Reset subsección al cambiar sección
      \fi
      % Verificar si necesitamos imprimir el título de subsección
      \ifx\leq@current@subsection\leq@last@printed@subsection\else
        \ifx\leq@current@subsection\@empty\else
          \addcontentsline{leq}{leqsubsection}{\leq@current@subsection}%
          \global\let\leq@last@printed@subsection\leq@current@subsection
        \fi
      \fi
      % Ecuación
      \refstepcounter{eqnum}%
      \begin{equation}
        #1\tag{E.\theeqnum}
      \end{equation}%
      \addcontentsline{leq}{leqt}{[E.\theeqnum] -- #2}%
      \addcontentsline{leq}{leqm}{#1}%
    }
    
    %%% ----- Índice de ecuaciones ----------%%%
    % Tamaño para []
    \newcommand{\leqIndexFont}{\normalsize}
    \newcommand{\leqTitleFont}{\tiny}
    \newcommand{\leqNoteFont}{\tiny}
    % Cabecera de sección 
    \newcommand\l@leqsection[2]{%
      \addvspace{25pt}% Mayor separación antes de nueva sección
      \noindent\leqIndexFont
      \makebox[\linewidth]{\rule{.38\linewidth}{0.4pt}\ \ \textbf{  \MakeUppercase{#1}}\ \ \rule{.38\linewidth}{0.4pt}}%
      \par\vspace{-10pt}
    }
    % Cabecera de subsección 
    \newcommand\l@leqsubsection[2]{%
      \par\addvspace{15pt}%
      \noindent\leqIndexFont #1
      \par\vspace{-7pt}
    }
    % Título de ecuación 
    \newcommand\l@leqt[2]{%
    \par\addvspace{10pt}%
      \par\noindent\leqTitleFont #1\par
    }
    % Miniatura de ecuación
    \newcommand\l@leqm[2]{%
      \vspace{0.3\baselineskip}%
      \par\scriptsize\noindent\hspace{1.5em}\(#1\)\par%
    }
    % Nota de ecuación 
    \newcommand\l@leqnote[2]{%
      \vspace{0.2\baselineskip}%
      \par\noindent\leqNoteFont\hspace{1.5em}\textbf{Nota.--} #1\par\vspace{0.5\baselineskip}%
    }
    % Generación del índice
    \newcommand{\mylistofequations}{%
      \clearpage
      \phantomsection
      % Título visual de la lista (ajústalo a tu gusto):
      {\centering\bfseries\Large Índice de ecuaciones\par}
      % Entrada en nuestro índice 'stc':
      \addcontentsline{stc}{section}{Índice de ecuaciones}%
      % Llamada a tu lista real (usa el nombre exacto de tu macro):
      \@starttoc{leq}%
    }
    
    %--------- Captura de títulos de sección --------%
    \let\leq@original@section\section
    \let\leq@original@subsection\subsection
    % Flag para saber si estamos en una sección asterisco
    \newif\ifLeqStarSection
    \LeqStarSectionfalse
    
    % Redefinir \section CON CONTROL DE ASTERISCO
    \renewcommand{\section}{%
      \@ifstar{\leq@section@star}{\@dblarg\leq@section@nostar}%
    }
    \def\leq@section@star#1{%
      \global\LeqStarSectiontrue
      \gdef\leq@current@section{#1}%
      \gdef\leq@current@subsection{}%
      \leq@original@section{#1}%
      \global\LeqStarSectionfalse
    }
    \def\leq@section@nostar[#1]#2{%
      \global\LeqStarSectionfalse
      \gdef\leq@current@section{#2}%
      \gdef\leq@current@subsection{}%
      \leq@original@section[#1]{#2}%
    }
    % Redefinir \subsection
    \renewcommand{\subsection}{%
      \@ifstar{\leq@subsection@star}{\@dblarg\leq@subsection@nostar}%
    }
    \def\leq@subsection@star#1{%
      \global\LeqStarSectiontrue
      \gdef\leq@current@subsection{#1}%
      \leq@original@subsection{#1}%
      \global\LeqStarSectionfalse
    }
    \def\leq@subsection@nostar[#1]#2{%
      \global\LeqStarSectionfalse
      \gdef\leq@current@subsection{#2}%
      \leq@original@subsection[#1]{#2}%
    }

    % ----- Restauración de valores Globales de estilo ----------%
    % Flags
    \newif\ifInFootnote
    \InFootnotefalse
    \pretocmd{\@footnotetext}{\InFootnotetrue}{}{}
    \apptocmd{\@footnotetext}{\InFootnotefalse}{}{}
    % Profundidad de listas (soporta anidadas)
    \newcount\MyListDepth \MyListDepth=0
    \AtBeginEnvironment{la}{\advance\MyListDepth by 1}
    \AfterEndEnvironment{la}{\advance\MyListDepth by -1}
    \AtBeginEnvironment{lnum}{\advance\MyListDepth by 1}
    \AfterEndEnvironment{lnum}{\advance\MyListDepth by -1}
    \AtBeginEnvironment{itemize}{\advance\MyListDepth by 1}
    \AfterEndEnvironment{itemize}{\advance\MyListDepth by -1}
    \AtBeginEnvironment{enumerate}{\advance\MyListDepth by 1}
    \AfterEndEnvironment{enumerate}{\advance\MyListDepth by -1}
    % Restauración diferida: hazlo al INICIO del próximo párrafo real
    \newif\if@pendingRestore
    \@pendingRestorefalse
    \newcommand{\RequestRestore}{\global\@pendingRestoretrue}
    \everypar{%
      \if@pendingRestore
        \global\@pendingRestorefalse
        \ifInFootnote\else
          \ifnum\MyListDepth=0
            \setlength{\parskip}{\GlobalParskip}%
            \setstretch{\GlobalBaseLineStretch}%
          \fi
        \fi
      \fi
    }
    % Al final de bloques:
    \AfterEndEnvironment{equation}{\RequestRestore}
    \AfterEndEnvironment{equation}{\RequestRestore}
    \AfterEndEnvironment{la}{\RequestRestore}
    \AfterEndEnvironment{lnum}{\RequestRestore}
    % Al final de macros:
    \apptocmd{\eqblock}{\RequestRestore}{}{}
    \apptocmd{\imagenfit}{\RequestRestore}{}{}
    
\makeatother

% ===================== ToC propio con 4 niveles (único bloque) =====================
\makeatletter

% --- Escritores: añadimos una línea al .stc en cada cabecera numerada ---
% Guardamos las versiones actuales para llamar al comportamiento normal
\let\MyTOC@orig@section\section
\let\MyTOC@orig@subsection\subsection
\let\MyTOC@orig@subsubsection\subsubsection
\let\MyTOC@orig@paragraph\paragraph

% SECTION
\renewcommand{\section}{%
  \@ifstar{\MyTOC@section@star}{\@dblarg\MyTOC@section@nostar}%
}
\def\MyTOC@section@star#1{%
  \MyTOC@orig@section{#1}% sin entrada al .stc
}
\def\MyTOC@section@nostar[#1]#2{%
  \MyTOC@orig@section[{#1}]{#2}% comportamiento normal (escribe al .toc)
  % Escribimos también nuestra línea al .stc (texto con \numberline)
  \addcontentsline{stc}{section}{\protect\numberline{\thesection}#2}%
}

% SUBSECTION
\renewcommand{\subsection}{%
  \@ifstar{\MyTOC@subsection@star}{\@dblarg\MyTOC@subsection@nostar}%
}
\def\MyTOC@subsection@star#1{%
  \MyTOC@orig@subsection{#1}%
}
\def\MyTOC@subsection@nostar[#1]#2{%
  \MyTOC@orig@subsection[{#1}]{#2}%
  \addcontentsline{stc}{subsection}{\protect\numberline{\thesubsection}#2}%
}

% SUBSUBSECTION
\renewcommand{\subsubsection}{%
  \@ifstar{\MyTOC@subsubsection@star}{\@dblarg\MyTOC@subsubsection@nostar}%
}
\def\MyTOC@subsubsection@star#1{%
  \MyTOC@orig@subsubsection{#1}%
}
\def\MyTOC@subsubsection@nostar[#1]#2{%
  \MyTOC@orig@subsubsection[{#1}]{#2}%
  \addcontentsline{stc}{subsubsection}{\protect\numberline{\thesubsubsection}#2}%
}

% PARAGRAPH
\renewcommand{\paragraph}{%
  \@ifstar{\MyTOC@paragraph@star}{\@dblarg\MyTOC@paragraph@nostar}%
}
\def\MyTOC@paragraph@star#1{%
  \MyTOC@orig@paragraph{#1}%
}
\def\MyTOC@paragraph@nostar[#1]#2{%
  \MyTOC@orig@paragraph[{#1}]{#2}%
  % Si no numeras los \paragraph, cambia \theparagraph por texto plano sin \numberline
  \addcontentsline{stc}{paragraph}{\protect\numberline{\theparagraph}#2}%
}

% --- Impresor del índice propio (lee el .stc) ---
\newcommand{\MyTOC}{%
  \par\bigskip
  {\centering\bfseries\Large Índice de contenido\par}%
  \medskip
  \begingroup
    \parindent=0pt\parskip=2pt
    \@starttoc{stc}%
  \endgroup
}

\makeatother
% ===================== fin ToC propio =====================


%%%%%%%%%%%%%%


% --- Datos del tema ---
\title{Política presupuestaria (I): Presupuestos del Estado\\
\large Marco legal e institucional de los presupuestos en España. Principios presupuestarios, con especial referencia a la estabilidad y sostenibilidad. Los Presupuestos Gnerales del Estado: elaboración, aprobación, ejecución y control. Clasificaciones de ingresos y gastos. Modificaciones presupuestarias}
\author{Marcelo Ortega}
\date{Septiembre 2026}

\begin{document}
\maketitle
\newpage

% --- Página de resumen y modificaciones ---
\puntosclave{ %% Escribir puntos clave Apuntes CECO
     Un reparto estimado del tiempo sería: (i) introducción, dos minutos; (ii) marco legal e institucional, cuatro minutos; (iii) principios presupuestarios,, cuatro minutos; (iv) ciclo presupuestario, trece minutos (elaboración ocho, aprobación uno, ejecución dos, control dos); (iv) clasificaciones de ingresos y gastos, cuatro minutos; (vi) modificaciones presupuestarias, dos minutos; (vii) conclusiones, uno minuto.
    }
\vspace{1cm}

\modificaciones{
    \textcolor{red}{\textbf{OJO ->} Indicaciones para los temas 21.B, 22.B y 23.B};

    Tema 4.B.21 -> Prestar el Marco jurídico-institucional de la actividad presupuestaria del Estado. Una vez presentado, debemos presentar las clasificaciones presupuestarias; esto es, la forma en que se presenta la actividad presupuestaria del estado atendiendo a criterios de \textbf{contabilidad pública} ([\authorfont{Ver Tema 4.B.22}]) y en términos de contabilidad nacional ([\authorfont{Ver Tema 4.B.23}]). 

    Tema 4.B.22 -> Todo lo presentado debe ir dirigido a presentar la actividad presupuestaria en términos de contabilidad presupuestaria.

    Tema 4.B.23 -> Todo lo presentado debe ir dirigido a presentar la actividad presupuestaria en términos de contabilidad nacional. 
    }

\newpage
\MyTOC
\newpage

\mylistofequations
\clearpage

\MyListOfFigures
\clearpage


\pagenumbering{arabic}
\setcounter{page}{1}


\section{Introducción}



\subsection{Contextualización}


\subsubsection{Relevancia}




\subsubsection{Problemática}



\subsection{Estructura de la exposición}


El presupuesto es el principal instrumento del sector público para implementar sus políticas. El presupuesto recoge una estimación de los ingresos que se prevén obtener en el t:iercicio presupuestario y el límite máximo de gastos que pueden realizarse.

El presupuesto es una herramienta de política económica que cumple múltiples funciones:
• Reflejo del plan económico del Gobierno. En la medida en qut' sus políticas impliquen algún gasto, el poder ejecutivo no podrá llevarlo a cabo a menos de que disponga de crédito suíiciente aprobado por una nonna con rango de ley: los presupuestos. El presupuesto no es un mero documento contable que registra de manera pasiva el plan económico del grupo político en el poder. sino que se constituye en una verdadera herramienta de política económica al decidir cómo distribuir unos recursos escasos entre unas necesidades casi infinitas.

• Herramienta de planificación y gestión. Además. el presupuesto pe1mite al sector público sistematizar la gestión, planificar sus actuaciones y evaluar su ejecución a través del seguimiento del gasto y de sus indicadores.

• Control de la labor del poder ejecutivo. Desde un punto de vista ex ante, el presupuesto limita las actividades que puede llevar a cabo el Gobierno. mientras que, desde una perspectiva ex post. pe1mite controlar el cumplimiento de los compromisos adquiridos en su proceso de elaboración y aprobación. Así, los presupuestos se podrían interpretar como un contrato entre la sociedad y su Gobierno (o, más concretameme. entre el legislativo yel ejecutivo).

La importancia no sólo económica sino también política de los presupuestos en España quedó patente en 2019, cuando el Gobierno central no fue capaz de aprobar en Cortes su presupuesto para ese año y decidió convocar nuevas elecciones.

Los Presupuestos Generales del Estado (en adelante PGE) se refieren, obviamente, al Estado. En contra de lo que pudiera parecer, el conjunto de los PGE consolidados gestiona poco más de la mitad del gasto público español, ya que el 44% es gestionado por Comunidades Autónomas (CCAA) y Entidades Locales (en adelante EELL). Ahora bien, los PGE condicionan las finanzas públicas del resto de Administraciones Territoriales (en adelante AATT). En efecto, los PGE determinan los recursos que van a recibir las CC AA y las EELL en vi11ud de los sistemas de financiación territoriales. En este sentido. hay que destacar que las transferencias que reciben las CCAA en virtud del sistema de financiación autonómico constituyen el ~80% de sus ingresos totales.

Como resumen podemos concluir que los PGE son:
• Una norma con rango de ley. 
• Una he,rnmienta de política fiscal y económica.
• Una herramienta de planificación y gestión.
• Una fuente de información de la actividad del sector público.
• Una vía de seguimiento y control de la actividad del poder ejecutivo.
• Un contrato entre la sociedad y su Gobierno (o, más concretamente. entre el legislativo y el ejecutivo).

En el primer apartado de este terna analizamos el marco legal e institucional en que se encuadra el ciclo presupuestario español. A continuación. analizaremos los principios presupuestarios, con especial énfasis en los principios de estabilidad presupuestaria y sostenibilidad financiera que quedaron reforzados con la reforma constitucional de 20 l l. El tercer apartado es el más importante de este tema. y desarrolla las distintas fases de los PGE. Después analizaremos cómo se clasifican los ingresos y los gastos públicos desde un punto de vista presupuestario. Por último. estudiaremos los distintos tipos de modificaciones presupuestarias que se pueden aplicar a unos PGE ya aprobados. Al final del tema se incluye un apartado de conclusiones.



\clearpage
\section{Actividad presupuestaria: Marco jurídico-conceptual}
\puntosclave{
    La regulación bá!>ica de los PGE se contempla en los artículos 1 34. 135 ) 136 de la Constitución Espa,iola. 
    
    El art. 1 34 regula. básicamente. la operativa presupuestaria. y ha sido desarrollado por la Ley General Presupuestaria. El art. 1 35, por su parte. fue modificado en 201 1 y recoge el principio de estabilidad presupuestaria. imponiendo el equilibrio presupuestario en términos estructurales y la prioridad del pago de la deuda y de sus i ntereses; este a,1 ículo ha sido desarrollado por la Ley Orgánica de Estabilidad Presupuestaria y Sostcnibilidad Financiera. El art. 1 36 regula el papel del Tribunal de Cuentas.
    
    A pesar de que lo!> PGE se aplican al sector público estatal. sus efectos trascienden al ámbito estatal. pues condicionan fuertemente los presupuestos de las CC/\A y de las EELL cuantificando sus sistemas de financiación y aprobando legislación de carácter básico que es de obligado cumplimiento para todas las Administraciones (p.ej. la revalorización de los salarios públicos).
}


Empecemos a presentar la actividad presupuestaria por donde cualquier institución jurídica merece: por su origen.

Antes de fijar qué dice hoy la Constitución sobre los PGE, o qué principios los rigen, conviene entender que ninguna de esas dos piezas nació de la nada en 1978. Por ello, este bloque recorre cuatro escalones: (i) el desarrollo histórico completo de la actividad presupuestaria en España; (ii) el marco jurídico actualmente vigente, situado ya con la perspectiva que la historia nos ha dado; (iii) el marco conceptual de principios que sostiene ese marco jurídico; y, (iv) el marco institucional del sector público presupuestario, con su doble desagregación horizontal y vertical.

\subsection{Desarrollo histórico de la actividad presupuestaria del Estado en España}
Empecemos por el principio. Antes de que existiese la palabra presupuesto en sentido técnico, existía el principio que la hace necesaria: ningún monarca medieval podía allegar recursos extraordinarios sin el consentimiento de quienes debían aportarlos. 

En los reinos medievales de Castilla y de Aragón, las Cortes otorgaban servicios a cambio de que el monarca atendiera determinadas peticiones o agravios de los procuradores, y era precisamente por el consentimiento a imponer esos subsidios extraordinarios (servicios) que las Cortes eran convocadas por el Rey. La misma lógica que el constitucionalismo inglés formularía después con precisión: \textit{no taxation without representation}.

Como vemos, en la Península ibérica esta lógica tiene una tradición propia y antigua, articulada en instituciones como las Cortes de León de 1188, frecuentemente señaladas como de las asambleas parlamentarias más tempranas de Europa con participación de representantes de las ciudades. Sin embargo, se fue erosionando a medida que la Corona conseguía, entre los siglos XVI y XVII, consolidar fuentes de ingreso permanentes y abundantes, ajenas al control efectivo de las Cortes. De este modo, para cuando llega el siglo XIX, España no hereda una tradición presupuestaria parlamentaria viva y continuada, sino un recuerdo medieval que el liberalismo gaditano tendrá que reinventar casi desde cero.

\subsubsection{Cádiz, 1811-1812: el nacimiento formal de la institución presupuestaria española}
El precedente directo de la institución presupuestaria moderna en España se sitúa en el Decreto de las Cortes de 22 de marzo de 1811, que inició la formación de presupuestos.

Como vemos, se origina incluso antes de que existiera el texto constitucional que les daría cobertura definitiva, que finalmente llegó con la Constitución de Cádiz, proumlgada en 1812. Este texto dedicó su Título VII (\textit{De las contribuciones}) a la Hacienda pública, con dieciocho preceptos que regulaban el establecimiento de las contribuciones por las Cortes, su periodicidad (anual), las clases de contribuciones, y su inclusión y distribución dentro del presupuesto. Además, dio cobertura orgánica a un órgano de control que resultará crucial: la Contaduría Mayor de Cuentas; antecedente directo del actual Tribunal de Cuentas, diferenciando ya entonces entre un control interno (Contadurías de valores y distribución) y un control de naturaleza jurisdiccional atribuido a la propia Contaduría Mayor.\footnote{%
    En 1819 se refundieron en esta Contaduría Mayor las llamadas contadurías especiales (valores, distribución y millones), que volverían a resurgir de forma separada en 1826, en el vaivén institucional típico de la Restauración absolutista fernandina, que alternó, según el momento político, entre restaurar y desmontar cada pieza del edificio liberal gaditano.
}

Tras la muerte de Fernando VII y el inicio de la Regencia de María Cristina, el Estatuto Real de 1834 trajo consigo una primera reforma de la contabilidad pública española, dentro del más amplio proceso de reconstrucción del Estado liberal.\footnote{%
    Es en este contexto donde, atendiendo a la literatura histórico-administrativa, el primer Presupuesto surgió en 1835. Momento en el que España cuenta con un documento reconocible de presupuesto general en sentido moderno, y no ya con la mera sucesión de rentas y gastos dispersos del Antiguo Régimen.
} No obstante, la pieza legislativa que consolidó y dio forma técnica a esta institución llegó quince años después, con la Ley de Contabilidad y Administración Pública de 1850, que instauró la contabilidad por ejercicios anuales. De esta forma, permitió por primera vez y de forma sistemática, la liquidación de presupuestos, la presentación de cuentas y el control efectivo por parte del Tribunal de Cuentas. Además, la Ley estableció que cada ministerio debía elaborar su propio presupuesto, génesis de la clasificación orgánica por secciones ministeriales que hoy sigue vertebrando los PGE, más de siglo y medio después.

### El problema que nunca se resolvió del todo: los créditos sin numerario
Conviene detenerse en un hilo que atraviesa la historia e ilustra un problema que trabajaremos con detalle más adelante: las reglas de gasto y disciplina fiscal.

Una crónica periodística española de principios del siglo XX, comentando la reforma de 1911, lo resume con una franqueza que sorprende. Denuncia: \textit{las autorizaciones que se consignan en todos los presupuestos para satisfacer diversas atenciones sin crédito numérico, atenciones más o menos imprevistas, que ponen muchas veces en peligro la nivelación del presupuesto; y, }[añade]\textit{ esta práctica la han prohibido ésta y todas las }[leyes]\textit{ que han regido desde 1850, unas veces en su letra y todas en su espíritu; lo que no ha sido obstáculo para usar y abusar de esos créditos, aplicados típicamente a las Clases Pasivas, que siempre se calculan deficientemente, al servicio de la Deuda, y a las ganancias de las Loterías y los monopolios estancados.}

Como vemos, este episodio ilustra la tentación de cualquier Hacienda de comprometer gasto por encima de lo presupuestado inicialmente, confiando en que la propia inercia administrativa termine legitimando el exceso. El mismo problema que hoy regulan con enorme precisión los créditos ampliables y los créditos extraordinarios; aunque sorprende comprobar que ya en 1850, y de nuevo en 1911, la norma trataba de cerrar exactamente esta misma puerta.

La Ley de Administración y Contabilidad de la Hacienda Pública, de 1 de julio de 1911, sustituyó a la ley de 1850 como marco regulador central de la actividad presupuestaria española, y su vigencia se prolongaría, con reformas puntuales, hasta prácticamente el final del franquismo.\footnote{%
    Por tanto, una longevidad institucional extraordinaria. Atravesó sin solución de continuidad la Restauración canovista, la Dictadura de Primo de Rivera, la Segunda República, la Guerra Civil y la práctica totalidad del régimen de Franco. Esta persistencia normativa tan prolongada contrasta con la enorme inestabilidad política del periodo que atraviesa. Además, ya hemos visto como esta época fue también de transformación profunda en el ámbito fiscal. Por tanto sorprende como, mientras el marco fiscal sustantivo español sufría transformaciones profundas, el armazón puramente procedimental de cómo se elabora, aprueba y controla el gasto público se mantuvo notablemente estable. Sugiriendo que las cuestiones de mecánica presupuestaria son menos sensibles a los vaivenes ideológicos que las cuestiones de a quién y cuánto se grava.
}

### De la Transición a la Constitución: la Ley General Presupuestaria de 1977

El final del franquismo trajo consigo una modernización de este marco procedimental: la Ley 11/1977, General Presupuestaria. Su aprobación, en enero de 1977, en el tramo final de la transición y bajo el Gobierno de la Monarquía que preparaba la apertura democrática, conviene leerla como parte del mismo impulso modernizador de la Hacienda española que, en el terreno tributario, culminaría poco después con la gran reforma fiscal de 1977-1978 liderada por FUENTES QUINTANA ([\authorfont{Ver Tema 4.B.12}]). Esta ley dotó a España de un marco presupuestario reconociblemente moderno, apenas dos años antes de que la Constitución de 1978 lo elevara a rango constitucional en sus arts. 134 a 136.

### De la Constitución al Texto Refundido de 1988, y de ahí a la Ley General Presupuestaria vigente

Tras la Constitución, la Ley 11/1977 fue objeto de sucesivas modificaciones que, con el paso de los años, hicieron aconsejable su refundición en un texto único. En 1988 se aprobó el Texto Refundido de la Ley General Presupuestaria, que rigió durante los años noventa y los primeros años dos mil, coincidiendo con el ingreso de España en la Comunidad Económica Europea, la adopción del euro, y las primeras exigencias de disciplina fiscal comunitaria. Finalmente, la Ley 47/2003, General Presupuestaria, hoy vigente, la sustituyó.

### Un cierre necesario: doscientos años para llegar al mismo problema de fondo

Si algo revela este recorrido de dos siglos, es que las grandes preguntas que rodean la actividad presupuestaria no son nuevas. Son, desde 1811, preguntas constitucionales sobre el reparto del poder entre quien gasta y quien controla el gasto; que no se inventaron en la Constitución de 1978, sino que heredó una tradición de doscientos años de intentos de responderlas con instrumentos jurídicos cada vez más sofisticados: quién controla el gasto público, con qué periodicidad, con qué límites frente al Ejecutivo, con qué mecanismo de rendición de cuentas.

\subsection{Marco jurídico actual de la actividad presupuestaria}




\subsubsection{Constitución Española}
La regulación básica de los PGE se contempla en los arts. 134-136 de la Constitución Española (CE).

El artículo 134 recoge algunos principios básicos de la actividad presupuestaria, así como aspectos del ciclo presupuestario que luego presentaremos. Por su parte, el reformado art. 135 gira en torno al principio de estabilidad presupuestaria y sus principales manifestaciones son el establecimiento de un límite al déficit estructural (las AAPP no podrán superar los límites desarrollados en LO, y las EELL deberán presentar equilibrio presupuestario),  un establecimiento al límites de deuda (con remisión expresa al TFUE) y la reclasificacion de los créditos por deuda como obligación prioritaria, sin que deuda o interés puedan ser objeto de enmienda o modificación alguna. Por último. el artículo 136 designa al Tribunal de Cuentas como el órgano de control de las cuentas y gestión económica de todo el sector público.\footnote{%
    Como sabemos, el Tribunal de Cuentas actual no es una institución de nuevo cuño diseñada en 1978. Pero además, su doble naturaleza actual, fiscalizadora (control de la actividad económico-financiera) y jurisdiccional (enjuiciamiento de las responsabilidades contables de quienes manejan fondos públicos), es la misma dualidad que el Decreto de 1813 ya establecía para su antecesor gaditano, diferenciando entonces entre un control interno y un control de naturaleza jurisdiccional. Esta combinación de funciones en un único órgano es, probablemente, la herencia institucional más directa y menos conocida de todo el proceso histórico que hemos recorrido.
}

La redacción exacta del articulado se expone a continuación:
\begin{lnum}
    \item \textbf{Art. 134 - Actividad presupuestaria} (i) Corresponde al Gobierno la elaboración de los Presupuestos Generales del Estado y a las Cortes Generales, su examen, enmienda y aprobación. (ii) Los Presupuestos Generales del Estado tendrán carácter anual, incluirán la totalidad de los gastos e ingresos del sector público estatal y en ellos se consignará el importe de los beneficios fiscales que afecten a los tributos del Estado. (iii) El Gobierno deberá presentar ante el Congreso de los Diputados los Presupuestos Generales del Estado al menos tres meses antes de la expiración de los del año anterior. (iv) Si la Ley de Presupuestos no se aprobara antes del primer día del ejercicio económico correspondiente, se considerarán automáticamente prorrogados los Presupuestos del ejercicio anterior hasta la aprobación de los nuevos. (v) Aprobados los Presupuestos Generales del Estado, el Gobierno podrá presentar proyectos de ley que impliquen aumento del gasto público o disminución de los ingresos correspondientes al mismo ejercicio presupuestario. (vi) Toda proposición o enmienda que suponga aumento de los créditos o disminución de los ingresos presupuestarios requerirá la conformidad del Gobierno para su tramitación. (vii) La Ley de Presupuestos no puede crear tributos. Podrá modificarlos cuando una ley tributaria sustantiva así lo prevea.
    \item \textbf{Art. 135 - Estabilidad presupuestaria}. (i) Todas las Administraciones Públicas adecuarán sus actuaciones al principio de estabilidad presupuestaria. (ii) El Estado y las Comunidades Autónomas no podrán incurrir en un déficit estructural que supere los márgenes establecidos, en su caso, por la Unión Europea para sus Estados Miembros. Una ley orgánica fijará el déficit estructural máximo permitido al Estado y a las Comunidades Autónomas, en relación con su producto interior bruto. Las Entidades Locales deberán presentar equilibrio presupuestario. (iii) El Estado y las Comunidades Autónomas habrán de estar autorizados por ley para emitir deuda pública o contraer crédito. Los créditos para satisfacer los intereses y el capital de la deuda pública de las Administraciones se entenderán siempre incluidos en el estado de gastos de sus presupuestos y su pago gozará de prioridad absoluta. Estos créditos no podrán ser objeto de enmienda o modificación, mientras se ajusten a las condiciones de la ley de emisión. El volumen de deuda pública del conjunto de las Administraciones Públicas en relación con el producto interior bruto del Estado no podrá superar el valor de referencia establecido en el Tratado de Funcionamiento de la Unión Europea. (iv) Los límites de déficit estructural y de volumen de deuda pública sólo podrán superarse en caso de catástrofes naturales, recesión económica o situaciones de emergencia extraordinaria que escapen al control del Estado y perjudiquen considerablemente la situación financiera o la sostenibilidad económica o social del Estado, apreciadas por la mayoría absoluta de los miembros del Congreso de los Diputados. (v) Una ley orgánica desarrollará los principios a que se refiere este artículo, así como la participación, en los procedimientos respectivos, de los órganos de coordinación institucional entre las Administraciones Públicas en materia de política fiscal y financiera. En todo caso, regulará: a) La distribución de los límites de déficit y de deuda entre las distintas Administraciones Públicas, los supuestos excepcionales de superación de los mismos y la forma y plazo de corrección de las desviaciones que sobre uno y otro pudieran producirse. b) La metodología y el procedimiento para el cálculo del déficit estructural. c) La responsabilidad de cada Administración Pública en caso de incumplimiento de los objetivos de estabilidad presupuestaria. (vi) Las Comunidades Autónomas, de acuerdo con sus respectivos Estatutos y dentro de los límites a que se refiere este artículo, adoptarán las disposiciones que procedan para la aplicación efectiva del principio de estabilidad en sus normas y decisiones presupuestarias.
    \item \textbf{Art. 136 - Tribunal de Cuentas}. El Tribunal de Cuentas es el supremo órgano fiscalizador de las cuentas y de la gestión económica de Estado, así como del sector público. (ii) Dependerá directamente de las Cortes Generales y ejercerá sus funciones por delegación de ellas en el examen y comprobación de la Cuenta General del Estado. (iii) Las cuentas del Estado y del sector público estatal se rendirán al Tribunal de Cuentas y serán censuradas por éste. (iv) El Tribunal de Cuentas, sin perjuicio de su propia jurisdicción, remitirá a las Cortes Generales un informe anual en el que, cuando proceda, comunicará las infracciones o responsabilidades en que, a su juicio, se hubiere incurrido. (v) Los miembros del Tribunal de Cuentas gozarán de la misma independencia e inamovilidad y estarán sometidos a las mismas incompatibilidades que los Jueces. (vi) Una ley orgánica regulará la composición, organización y funciones del Tribunal de Cuentas.
\end{lnum}



\paragraph{Contenido: ¿qué puede contener la Ley de PGE?}
Hay una previsión concreta del art. $135$ que merece un desarrollo propio, porque ha generado uno de los debates más persistentes del Derecho constitucional. Como vemos, la Constitución prohíbe expresamente que la Ley de Presupuestos cree tributos, permitiéndole únicamente modificarlos cuando una ley tributaria sustantiva así lo prevea. 

No obstante, el Tribunal Constitucional ha ido más allá limitando también qué otras materias, más allá de lo estrictamente tributario, pueden colarse en el articulado de la Ley de Presupuestos. Según doctrina reiterada (STC 248/2007), solo pueden regularse en ella materias que \textbf{guarden una relación directa e inmediata con los ingresos o gastos del Estado, respondan a los criterios de política económica del Gobierno, o sirvan a una mayor inteligencia o mejor ejecución del presupuesto}.

Esta limitación constitucional explica un fenómeno que conviene conocer, porque ilustra cómo una restricción jurídica puede generar un problema colateral de técnica normativa. Se trata de las llamadas \textit{leyes de acompañamiento}, normas ordinarias tramitadas en paralelo a la Ley de Presupuestos de cada ejercicio, y aprobadas casi simultáneamente con ella, que buscan dar cobijo a todas aquellas medidas de política económica y social que el Gobierno quiere sacar adelante pero que la propia jurisprudencia constitucional le impide introducir en la Ley de Presupuestos por carecer de esa \textbf{relación directa e inmediata}. 

La crítica doctrinal a esta práctica señala una paradoja evidente. El límite al contenido de la Ley de Presupuestos, pensado para proteger la calidad del debate parlamentario sobre materias tributarias y evitar que se cuelen otras cuestiones aprovechando la tramitación urgente y prioritaria del presupuesto, termina generando una segunda norma que en la práctica reproduce el mismo problema que se pretendía evitar. Un texto heterogéneo, de contenido muy disperso, tramitado con la misma urgencia de fin de año y con escaso margen real de debate sobre sus disposiciones.

\paragraph{Reforma: ¿cómo se hizo, merece ser criticada?}
Hemos presentado el contenido sustantivo del art. 135 CE reformado: déficit estructural, límite de deuda por remisión y prioridad absoluta de pago de la deuda. Lo que conviene desarrollar es el propio episodio de la reforma, porque es uno de los momentos más singulares de la historia constitucional reciente.

La Proposición de reforma fue presentada conjuntamente por los grupos parlamentarios Socialista y Popular el 26 de agosto de 2011, en pleno periodo estival y con un Parlamento que ya tenía fecha de caducidad. Las elecciones generales estaban convocadas para el 20 de noviembre de ese mismo año. Los dos grupos mayoritarios solicitaron su tramitación por el procedimiento de urgencia y lectura única, evitando así el debate en Comisión artículo por artículo. El Congreso la aprobó el 2 de septiembre de 2011; el Senado, sin introducir ninguna modificación al texto remitido, el 7 de septiembre; ambas cámaras alcanzaron la mayoría reforzada de tres quintos que exige el art. 167 CE para la reforma constitucional por vía ordinaria. 

Transcurridos los quince días que ese mismo artículo concede para que una décima parte de cualquiera de las cámaras solicite someter la reforma a referéndum, grupos como Izquierda Unida-Iniciativa per Catalunya Verds y el BNG lo intentaron activamente, presentando enmiendas a la totalidad que reclamaban acudir a la vía agravada del art. 168 CE, con referéndum vinculante obligatorio. Sin embargo, esa solicitud no prosperó, y el texto se publicó en el BOE el 27 de septiembre de 2011. En total, trascurrieron trece días desde la presentación de la proposición hasta su aprobación definitiva en el Senado. La segunda reforma en toda su historia, tras la modificación, mucho más técnica y consensuada, del art. $13.2$ en 1992 para adaptar el derecho de sufragio pasivo a las exigencias del Tratado de Maastricht.

Como es normal, el procedimiento generó una resentimiento notable. La catedrática de Derecho Constitucional GARCÍA-ESCUDERO MÁRQUEZ la analizó con detalle en dos trabajos casi coetáneos, según la cual los plazos reducidos y el procedimiento de lectura única no resultaban adecuados para una reforma constitucional, precisamente por impedir el sosiego y el intento de consenso con los grupos minoritarios que este tipo de reformas debería exigir. Junto a esta crítica procedimental, el debate de fondo que la reforma suscitó. Los críticos de la reforma sostuvieron, por un lado, que era jurídicamente innecesaria, porque los mismos límites de déficit y deuda ya estaban establecidos, con suficiente detalle, tanto en la legislación española de estabilidad presupuestaria vigente desde mediados de los 2000 como en el propio Derecho de la Unión Europea, de aplicación directa desde 1997. Y, por otro, que su verdadero efecto era constitucionalizar una determinada orientación de Política Económica (disciplina fiscal como prioridad absoluta), un terreno que hasta entonces se había considerado propio del debate político ordinario y no de la arquitectura constitucional básica. Además, algunos señalaron otro efecto colateral sobre la autonomía financiera de Comunidades Autónomas y Corporaciones Locales. 

Los defensores de la reforma, por su parte, la justificaron por la necesidad de enviar una señal de credibilidad reforzada a los mercados financieros y a las instituciones europeas en el momento álgido de la crisis de deuda soberana, en un contexto en el que Alemania ya había incorporado a su propia Ley Fundamental, en 2009, un freno constitucional al endeudamiento (\textit{Schuldenbremse}) que sirvió de referencia explícita al debate español, y que se adelantaba, además, a lo que meses después sería el propio Tratado de Estabilidad, Coordinación y Gobernanza de la Unión (el Pacto Fiscal), que exigiría a sus firmantes incorporar reglas de equilibrio presupuestario de rango preferentemente constitucional.

\subsubsection{Legislación de desarrollo: ¿cuál es su cuerpo normativo?}
En un rango normativo inferior a la Constitución española se sitúan las leyes que regulan la materia presupuestaria (respetando siempre los preceptos constitucionales).

Salvando generalidades, podríamos decir que el art. $134$ se desarrolla mediante la Ley 47/2003 General Presupuestaria (LGP), mientras que el art. $135$ se desarrolla por Ley Orgánica 2/2012 (LOEP).

Sobre la relación entre ambas conviene fijar que no existe jerarquía, por más que la ley orgánica exija una mayoría cualificada para su aprobación en las Cortes. El Tribunal Constitucional ha determinado, en numerosas ocasiones, que la relación entre ambos instrumentos es de materia. La Constitución reserva a la ley orgánica el desarrollo legal de determinadas materias tasadas, de modo que esas materias no pueden legislarse mediante ley ordinaria. Y viceversa, una ley orgánica no puede invadir, sin habilitación constitucional expresa, el terreno reservado a la ley ordinaria. Esta arquitectura explica por qué la propia LOEP dedica buena parte de su articulado a cuestiones de detalle técnico-contable (cálculo del déficit estructural, metodología de la regla de gasto) que, en un sistema distinto, bien podrían haberse dejado a un reglamento de desarrollo. La propia naturaleza del art. 135 CE, al exigir explícitamente que sea una ley orgánica la que desarrolle sus límites concretos, obliga a elevar a ese rango reforzado un nivel de detalle técnico que en la mayoría de ordenamientos comparados quedaría en sede reglamentaria, prolongando así la misma lógica de constitucionalización o refuerzo normativo.

Otra diferencia importante es que, mientras que la LGP se desarrolla en términos de contabilidad pública o presupuestaria (más próxima al principio de caja), la LOEP se desarrolla en términos de contabilidad nacional; es decir, el Sistema Europeo de Cuentas, que representa el devengo puro.

\paragraph{Ley General Presupuestaria (LGP)}
La Ley 47/2003, General Presupuestaria (LGP), que sabemos desarrolla el art. $134$ CE, regula el régimen presupuestario, económico-financiero, de contabilidad, de intervención y de control financiero.  El objeto central del tema es su vertiente sobre el régimen presupuestario que regula, el alcance y contenido de los PGE, así como su régimen procedimental. 

El alcance subjetivo de la LGP es el sector público estatal, que pueden dividirse en entidades con presupuesto limitativo  (no pueden sobrepasar el gasto presupuestado) y las entidades de presupuesto estimativo (pueden sobrepasar el gasto presupuestado). Las entidades con presupuesto limititativo son las que forman parte del presupuesto consolidado del Estado, que se forman agregando ingresos y gastos de estas entidades, y restando las transferencias internas entre ellas para evitar la doble contabilidad.

Como es natural, si lo ponemos en contexto, a pesar de que el ámbito subjetivo se circunscriba al sector público \textbf{estatal}, sus efectos trascienden enormemente a otros niveles de gobierno por dos vías. Por un lado, los PGE determinan en gran medida los ingresos de las CCAA y EELL, ya que en el contexto de su elaboración el Gobierno cuantifica los sistemas de financiación autonómico y local.\footnote{
En la elaboración de los presupuestos quedan determinados los sistemas de financiación autonómico y local. En lo relativo al sistema de autonómico, durante su elaboración el Gobierno determina los dos pilares del sistema. Por un lado, las cuantías de los fondos del sistema. Por otro, las previsiones de ingresos tributarios.

\textbf{NOTA.-} Las cuotas correspondientes a ingresos tributarios participados y/o cedidos no figuran en los PGE, por lo que la cifra de los PGE se refiere únicamente a la recaudación que se queda el Estado. No obstante, en la mayoría de impuestos cedidos, los cobra íntegramente el Estado y luego transfiera a las CCAA su participación (algo que no ocurre con los impuestos cedidos, que los recaudan directamente las CCAA). Por el contrario, los Fondos del Sistema sí figuran en los PGE. Como gastos se contabilizan las entregas a cuenta y liquidaciones que hace el Estado, y como ingresos las entregas a cuenta y liquidaciones de las CCAA que les sale a pagar.

También figuran en los PGE los recursos que el Estado transfiere a las CCAA por convenios, contratos-programa o subvenciones para financiar determinados bienes o servicios prestados por las CCAA. Por ejemplo, políticas activas de empleo, dependencia, proyectos de inversión relacionados con el Plan de Recuperación, Transformación y Resiliencia. etc.
} Por otro lado, los PGE también condicionan sus gastos ya que su articulado suele contener disposiciones de carácter básico y, en consecuencia, obligado cumplimiento para todas las Administraciones (por ejemplo, la revalorización de los salarios públicos).

\paragraph{Ley Órganica de Estabilidad Presupuestaria (LOEP)}
Por otro lado, la Ley Orgánica de Estabilidad Presupuestaria y Sostenibilidad Financiera de 2012 (LOEP) regula los principios, límites, procedimientos e instrumentos que permiten asegurar la estabilidad presupuestaria.

Su ámbito de aplicación subjetivo se extiende a todas las Administraciones Públicas, a diferencia de la LGP. Es decir, incluye tanto al Estado y Seguridad Social, como a las CCAA y EELL.

\subsection{Marco conceptual de la actividad presupuestaria: los principios}
Fijado ya el marco jurídico, conviene preguntarse por qué ese marco dice lo que dice: ¿qué principios sostienen, con qué justificación económica y política, cada una de las reglas que hemos ido enumerando?

La LOEP amplió los principios que ya preveía la LGP, dividiéndolos en dos familias: los de programación presupuestaria y los de gestión presupuestaria. Los primeros tienen una vocación más general (LOEP), de planificación, mientras que los segundos se refieren más a los principios que rigen el ciclo presupuestario (LGP).

\subsubsection{LOEP: Principios de programación presupuestaria}
Los principios de programación presupuestaria son siete y se regulan en los arts. $135$ de la CE, $3$ a $9$ de la LOEP y $26$ de la LGP.

\paragraph{Estabilidad presupuestaria: ¿qué indicador adoptar?}
El principio de estabilidad presupuestaria, atendiendo a la LOEP, sería la situación de equilibrio o superávit estructural, prohibiendo el déficit estructural salvo dos excepciones: (i) reformas estructurales con efectos a largo plazo, hasta el $0,4\%$ del PIB; y, (ii) situaciones excepcionales, con cláusula de escape apreciada por mayoría absoluta del Congreso.\footnote{%
    En el caso de España, a petición del Gobierno, el Congreso apreció situación de emergencia en los años 2020-2023. La apreciación de la cláusula de escape supone la no sujeción a objetivos fiscales durante esos años, aunque exige la aprobación posterior de un plan de reequilibrio.
}

La pregunta que conviene hacerse es por qué el legislador eligió el déficit estructural como magnitud de referencia, en lugar del déficit nominal, mucho más sencillo de observar. Como sabemos, el déficit nominal se mueve por razones ajenas a cualquier decisión de política fiscal deliberada, de modo que exigir el mismo esfuerzo de ajuste sobre el déficit nominal en una recesión que en una expansión obligaría a recortar gasto precisamente cuando la economía más lo necesita, agravando el propio ciclo (problema clásico de prociclicidad). Por el contrario, el déficit estructural pretende filtrar ese componente cíclico, aislando la parte del déficit que sí depende de decisiones discrecionales del Gobierno. Sin embargo, esto esconde un problema práctico de primer orden. El déficit estructural no se observa directamente, sino que se estima.\footnote{%
    Se define como el déficit ajustado del ciclo neto de medidas excepcionales y temporales. Lo calcula el Ministerio de Economía siguiendo la metodología europea, sobre el cálculo de las previsiones tendenciales de ingresos y gastos y de la tasa de referencia de la economía española.
} Como es natural, esa estimación depende de una metodología con supuestos técnicos discutibles y con una notoria inestabilidad. La propia Comisión Europea ha revisado sus propias estimaciones de output gap y de saldo estructural de años ya cerrados, cuando en su momento esas mismas cifras habían servido de base para exigir, o para eximir, un ajuste fiscal concreto a un Estado miembro. En el fondo, exigir el ajuste sobre una magnitud teóricamente superior, pero que cambia de valor según quién y cuándo la calcule, genera un problema de seguridad jurídica y de previsibilidad para cualquier Gobierno que intente cumplirla de buena fe.

Por último, de este principio deriva la \textbf{regla de gasto} (art. $12$ LOEP). El gasto computable de la Administración Central, de las CCAA y de las EELL no puede crecer por encima de la \textit{tasa de referencia} (el crecimiento del PIB nominal español a medio plazo).\footnote{%
    El gasto computable se define como el gasto total de la Administración en cuestión, detrayendo una serie de cuantías: medidas de ingresos que aumenten su recaudación de manera permanente, intereses de la deuda, gasto cíclico en prestaciones por desempleo, gasto financiado con fondos europeos, y transferencias a CCAA y EELL asociadas a los sistemas de financiación.

\textbf{NOTA.-} Esta regla de gasto es muy similar a la prevista en el PEC (expenditure benchmark), aunque no es del todo idéntica, pues afecta a todos los subsectores sin excluir a la Seguridad Social, resta los oneoffs y suaviza el gasto en inversión, entre otras divergencias.
} Conviene señalar, que esta regla de gasto ha ido ganando peso relativo dentro de la propia gobernanza fiscal europea, hasta convertirse en el instrumento operativo central de la reforma del marco fiscal de la UE que entró en vigor en 2024: la razón es que una regla basada en el gasto resulta menos vulnerable a los shocks de ingresos y, por tanto, menos susceptible de sesgo procíclico que las reglas basadas directamente en el saldo, precisamente el problema que motivó la propia elección del criterio estructural frente al nominal.

\paragraph{Sostenibilidad financiera: ¿qué criterio elegir?}
El principio de sostenibilidad financiera exige, atendiendo a la LOEP, la capacidad para financiar compromisos de gasto presentes y futuros dentro de los límites de déficit, deuda pública y morosidad de deuda comercial. ¿Cómo? A priori, limitando la deuda al $60\%$ del PIB para el conjunto de AAPP ($44\%$ AACC, $13\%$ CCAA, $3\%$ EELL). 

Ahora bien, como podemos intuir, este nivel es una creación política más que un óptimo riguroso. El $60\%$ procede del Tratado de Maastricht de 1992, y su origen no es el resultado de ningún cálculo económico sobre cuál sería el nivel óptimo de deuda pública; sino, simplemente, la deuda pública media de los entonces doce Estados miembros de la Comunidad Europea a finales de los años ochenta y principios de los noventa, tomada como techo razonable sin que exista ninguna demostración rigurosa de que ese nivel maximice el bienestar de los ciudadanos europeos o de que constituya un óptimo teórico.\footnote{
El hecho de que las Cortes decidiesen incorporar un criterio heredado, sin más, a la Constitución, deja mucho que desear. Además, se entiende con mayor razón las críticas que se vertieron a esta reforma \textit{exprés} del texto constitucional.
}

El propio $3\%$ de déficit, su compañero inseparable, guarda con el $60\%$ una relación aritmética que TIETMEYER (negociador alemán de Maastricht) explicitó después: si se asume un crecimiento nominal del PIB del $5\%$ ($3\%$ real más $2\%$ de inflación), una deuda estabilizada en el $60\%$ del PIB es exactamente compatible con un déficit anual constante del $3\%$. Es decir: el $3\%$ no se calculó de forma independiente; se calculó para ser coherente con el $60\%$, bajo unos supuestos de crecimiento que hoy, con tipos de interés estructuralmente más bajos y un crecimiento potencial menor, ya no se sostienen. De hecho, el propio debate entorno a la reforma del marco fiscal europeo parte de reconocer que los cambios estructurales de los últimos años sitúan a las economías europeas en un escenario completamente diferente al de 1992.

\begin{specialblock}{\textbf{Reglas fiscales}: Diferencias}
    Las diferencias entre las reglas de déficit y deuda y la regla de gasto son:
    \begin{la}
        \item \textbf{Ámbito subjetivo}. Las reglas de déficit y deuda afectan a todos los subsectores de las AAPP, mientras que la regla de gasto afecta a todos los subsectores menos a la Seguridad Social;
        \item \textbf{Automaticidad}. Las reglas de déficit y deuda tienen que ser aprobadas por las Cortes. La regla de gasto es de aplicación automática;
        \item \textbf{Diferenciación}. Las reglas de déficit y deuda se fijan para el total de las AAPP, así como para cada subsectore. La regla de gasto es única para el conjunto.
    \end{la}
\end{specialblock}

\paragraph{Plurianualidad: ¿qué prevalece, anualidad o consistencia temporal?}
La elaboración de los presupuestos debe encuadrarse en una visión de medio plazo. La justificación de fondo conecta con un problema de economía política bien documentado. 

Cualquier Gobierno, sometido a un ciclo electoral relativamente corto, tiene un incentivo estructural a descontar excesivamente el futuro frente al presente. El mismo problema de inconsistencia temporal que la literatura de economía política (ALESINA, TABELLINI y otros) ha documentado repetidamente en democracias comparadas. Un marco de planificación plurianual, que obliga a fijar compromisos de déficit y deuda para los tres años siguientes y no solo para el ejercicio en curso, pretende introducir un contrapeso institucional frente a esa miopía electoral. Pero merece señalarse la tensión genuina que este principio mantiene con el de anualidad: cuanto más se alarga el horizonte de planificación, más difícil resulta también el control parlamentario efectivo. Las Cortes aprueban el presupuesto de un año, pero solo ratifican objetivos agregados de déficit y deuda para los tres siguientes, sin la misma capacidad de escrutinio. Plurianualidad y anualidad no son, por tanto, dos principios perfectamente complementarios. Cada uno resuelve un problema que el otro, por sí solo, no puede resolver, y su convivencia es, en cierto sentido, una tensión nunca cerrada.

\paragraph{Transparencia: ¿exigencia democrática?}
Las cuentas de las AAPP deben contener información suficiente para verificar su situación financiera y el cumplimiento de los objetivos fiscales.

El Open Budget Survey, elaborado periódicamente por el International Budget Partnership, evalúa la transparencia presupuestaria de decenas de países según criterios homogéneos de disponibilidad, oportunidad y detalle de la documentación presupuestaria pública. España obtiene, en las sucesivas ediciones de este índice, puntuaciones que la sitúan en el grupo intermedio-alto de los países evaluados, muy lejos de los líderes escandinavos, pero también lejos de los países con menor transparencia.

\paragraph{Eficiencia: ¿el criterio más controvertido de todos?}
La asignación y utilización de los recursos públicos debe someterse al principio de eficiencia. La idea no nace con la LOEP, tiene una genealogía propia que se remonta al movimiento de reforma presupuestaria del Progresismo estadounidense de comienzos del siglo XX.\footnote{%
    La Comisión Taft sobre Economía y Eficiencia (1910-1912), encargada por el presidente Taft, formuló por primera vez de forma sistemática la idea del executive budget: un presupuesto único, elaborado con criterios técnicos y no como mera agregación de peticiones departamentales inconexas. Bajo la influencia directa de las ideas de gestión científica de TAYLOR que por entonces revolucionaban la organización industrial privada, autores como CLEVELAND (1915) y WILLOUGHBY (1918) sistematizaron esta doctrina, que cristalizaría en la Budget and Accounting Act de 1921, creadora de la moderna Oficina de Presupuesto de EEUU.
}

La primera implementación técnica más ambiciosa de este principio llegó con el PPBS (Planning-Programming-Budgeting System). Introducido por McNamara en el Departamento de Defensa estadounidense a comienzos de los años sesenta, se extendió poco después, por orden del presidente Johnson, a la totalidad del Gobierno federal. Es el antecedente directo de nuestra propia presupuestación por programas. 

Ahora bien, respecto a este principio se vertió una de las críticas más influyentes de la ciencia presupuestaria. WILDAVSKY (1964) sostuvo que la presupuestación real, a diferencia de la que prescribe la teoría racional-comprensiva del PPBS o el principio de eficiencia, es incrementalista: los actores presupuestarios no reevalúan cada año, desde cero, la totalidad del gasto de cada política a la luz de sus fines últimos, sino que negocian sobre una base ya heredada del año anterior, discutiendo solo el incremento o decremento marginal, porque la negociación política hace inviable, en la práctica, la reevaluación total y racional que exigía.\footnote{%
    El propio Gobierno federal estadounidense abandonó formalmente el PPBS como sistema de aplicación obligatoria en 1971, reconociendo implícitamente el diagnóstico de WILDAVSKY.
}

Desde esta óptica, la propia AIReF ha constatado que la presupuestación por programas, pensada para introducir precisamente este enfoque de eficiencia y resultados, tiene una aplicación más formal que real. El diagnóstico no es exclusivo de España, sino que es la misma crítica que la literatura de administración pública dirige al movimiento de la Nueva Gestión Pública (New Public Management). La introducción de indicadores, objetivos y evaluación de resultados en el sector público resulta, con notable frecuencia, más sencilla de aprobar formalmente en la norma que de aplicar, porque exige una capacidad técnica de medición y una cultura organizativa de rendición de cuentas que no se generan aprobando una ley.

La corriente alternativa más reciente, los \textit{spending reviews} (que AIReF realiza desde 2017), representa un intento más modesto y más viable de recuperar el espíritu del principio sin repetir la ambición totalizadora, y el consiguiente fracaso, del planteamiento original.

\paragraph{Responsabilidad: ¿son creíbles los mecanismos sancionadores?}
Atendiendo al principio de responsabilidad, las AAPP que incumplan los objetivos fiscales asumirán la parte imputable de las responsabilidades derivadas. 

Esta exigencia tiene su fundamento teórico en la literatura macroeconómica inaugurada por KYDLAND y PRESCOTT (1977), trabajo que les valdría el Premio Nobel de Economía. Su argumento central, extendido con naturalidad a la política fiscal, es el de la inconsistencia temporal. Un gobierno que anuncia hoy una senda de disciplina fiscal futura tiene, en el momento de cumplir esa promesa, un incentivo a desviarse, de modo que la promesa, por sí sola y sin ningún mecanismo de cumplimiento forzoso, carece de credibilidad ante los agentes económicos. Por tanto, un régimen de responsabilidad con consecuencias reales es la única forma de hacer creíble un compromiso.

Conviene señalar, que este principio comparte una debilidad bien documentada con su equivalente europeo. El Pacto de Estabilidad y Crecimiento, precursor directo de esta idea, ha sido objeto de crítica recurrente por su falta de credibilidad sancionadora. En este mismo sentido, la AIReF ha señalado que las medidas preventivas, correctivas y coercitivas de la LOEP tampoco se aplican siempre en el ámbito interno español; presentando el mismo problema de fondo a escala doméstica: la solidez teórica no garantiza su aplicación efectiva cuando el incumplidor tiene peso político suficiente para negociar su propio incumplimiento.

Existe también otra crítica desde otro frente, de naturaleza distinta. La propia literatura posterior, con economistas como BLANCHARD, ha señalado que un régimen de reglas excesivamente rígido, diseñado precisamente para resolver el problema de la inconsistencia temporal, puede terminar generando el problema inverso: una rigidez excesiva que impide al Gobierno responder con la flexibilidad necesaria ante una crisis genuina, precisamente la tensión que las cláusulas de escape de la LOEP (activadas en España 2020-2023) intentan resolver. Cuanto más fácil resulta activar la cláusula de escape, menos creíble es la regla en tiempos normales; cuanto más difícil, menos capacidad de respuesta tiene el Gobierno en una emergencia real. 

\paragraph{Lealtad institucional: ¿pueden internalizarse las externalidades?}
Conforme al principio de lealtad institucional, cada Administración debe valorar el impacto de sus actuaciones sobre las demás, respetar el reparto competencial y prestar cooperación. Este principio traduce un problema clásico de la teoría del federalismo fiscal, OATES (1972): en cualquier sistema de gobierno descentralizado, las decisiones de gasto o de endeudamiento de un nivel de gobierno generan externalidades sobre los demás niveles, de modo que ninguna regla fiscal agregada tiene sentido sin mecanismos explícitos de coordinación entre ellos. 

La aplicación más precisa de este marco en sistemas descentralizados la formula RODDEN (2006). El autor documenta cómo, en ausencia de mecanismos creíbles de coordinación, los gobiernos subnacionales de un Estado descentralizado tienden a desarrollar una expectativa racional de rescate por parte del gobierno central en caso de dificultades financieras (restricción presupuestaria blanda), generando un incentivo sistemático al sobreendeudamiento subnacional. 

El propio Fondo de Liquidez Autonómica español y la condicionalidad que impone a las Comunidades Autónomas que acceden a él, es la aplicación práctica más directa de esta literatura. Un mecanismo de rescate real, pero condicionado, diseñado precisamente para que la expectativa de rescate no se traduzca en un incentivo sistemático a la indisciplina previa.


\subsubsection{Principios de gestión presupuestaria}
Antes de entrar en cada uno, conviene fijar su origen común. En su práctica totalidad, son la sistematización que la doctrina clásica europea de finanzas públicas hizo, entre finales del siglo XIX y comienzos del XX, de los principios que el constitucionalismo liberal había ido decantando desde Cádiz. 

El jurista francés JÈZE () es probablemente el sistematizador más citado de estos principios clásicos del presupuesto, recogiendo y ordenando una cuerpo doctrinal que incluye a LEROY-BEAULIEU (1877). Esta doctrina continental fue recibida y adaptada por la ciencia jurídico-financiera española del siglo XX hasta llegar a la formulación actual del art. 27 LGP.

Actualmente, son seis principios y se regulan en los arts. $134$ de la CE y $27$ de la LGP.

\paragraph{Competencia: ¿qué puede realmente hacer cada institución?}
El principio de competencia exige que la elaboración y ejecución correspondan al Gobierno, la aprobación a las Cortes, y el control a los tres poderes.

Es, sin necesidad de forzar la analogía, la aplicación directa del principio de separación de poderes de MONTESQUIEU al terreno específico del gasto público: ningún poder controla, en solitario, el ciclo completo del dinero público, precisamente para que ninguno pueda gastar sin someterse al escrutinio de otro. 

Relajación: modificaciones presupuestarias (parte de ellas sin participación de las Cortes) y la prerrogativa del Gobierno de vetar enmiendas que aumenten gasto o disminuyan ingresos. 

No es una simple aplicación genérica de la separación de poderes. Tiene una genealogía constitucional propia y muy concreta, centrada específicamente en el poder de la bolsa (power of the purse), cuyo hito fundacional es la Declaración de Derechos inglesa de 1689 (Bill of Rights), que consagró el monopolio del Parlamento sobre la aprobación de impuestos y, progresivamente a lo largo del siglo XVIII, sobre la autorización del gasto público. 

\begin{specialblock}{\textbf{BISMARCK vs. Landtag:} Ejemplo histórico de la competencia}
    El episodio que mejor ilustra qué ocurre cuando este principio de competencia se rompe realmente es el conflicto constitucional prusiano de 1862-1866. Cuando el Landtag liberal se negó a aprobar el presupuesto militar reclamado por el rey Guillermo I, el recién nombrado ministro-presidente Otto von Bismarck recurrió a la llamada Lückentheorie (teoría de la laguna), formulada por STAHL:\footnote{%
        Como anécdota, BISMARCK fue nombrado el 23 de septiembre de 1862 y, poco después pronunciaría precisamente ante la comisión de presupuestos del Parlamento, su célebre discurso de \textit{hierro y sangre}: 

    [...] \textit{Los grandes problemas de la época no se resolverán con discursos y decisiones tomadas por mayoría —éste fue el tremendo error de 1848 y 1849—, sino con sangre y hierro}
    } ante la falta de acuerdo entre Corona y Parlamento sobre el presupuesto, sostuvo Bismarck, existía una laguna constitucional que el principio monárquico estaba legitimado para colmar, permitiendo al Gobierno recaudar impuestos y realizar gastos sin presupuesto aprobado. Prusia funcionó así durante cuatro años completos, hasta que la victoria militar prusiana sobre Austria en Königgrätz (1866) debilitó políticamente a la oposición liberal lo bastante como para que aceptara una Ley de Indemnidad que legalizó retroactivamente todos los presupuestos ejecutados sin aprobación entre 1862 y 1866. No es, por tanto, una disputa abstracta sobre teoría constitucional, sino cuatro años de un Estado gastando y recaudando sin el consentimiento parlamentario que la propia Constitución prusiana exigía, resuelto finalmente no por la fuerza del argumento jurídico sino por la fuerza bruta de una victoria militar que cambió la correlación de fuerzas políticas.
\end{specialblock}

\paragraph{Anualidad}
El principio de anualidad exige que los gastos e ingresos contenidos en los PGE se refieren al año natural, que se extiende del 1 de enero al 31 de diciembre. El principio de anualidad es un tanto particular y, como adelantamos, se encuentra en tensión con la plurianualidad.

Encuadrada dentro de la teoría política liberal, MILL (1861), sitúa el control periódico y renovado del gasto público entre las funciones centrales de cualquier asamblea representativa, precisamente porque un consentimiento que no caduca y no exige renovación periódica deja de operar, con el paso del tiempo, como un control real. Se convierte, argumenta MILL, en una autorización cada vez más formal y cada vez menos sustantiva: el Parlamento no revisa lo ya aprobado, y el Gobierno no siente la necesidad de justificarlo de nuevo. Es, con ciertas adaptaciones, el mismo argumento de KYDLAND y PRESCOTT: un compromiso sin renovación periódica pierde, con el tiempo, su capacidad real de vincular a quien lo asumió.

\paragraph{Unidad y universalidad}
Atendiendo al principio de unidad y universalidad, los presupuestos deben recoger todos los ingresos y gastos del sector público estatal. 

Es, junto con la anualidad, el principio más directamente jezeiano de todos. Para JÈZE, un presupuesto fragmentado en múltiples documentos y fondos distintos, sin una visión de conjunto única, impide al Parlamento ejercer un control real sobre el volumen total de la actividad financiera del Estado, porque ningún control puede ser completo si no conoce la totalidad de lo que pretende controlar.

La literatura moderna de transparencia fiscal ha documentado extensamente el riesgo concreto de incumplir dicho principio: la proliferación de fondos extrapresupuestarios, como técnica deliberada para sustraer determinadas partidas de gasto del escrutinio parlamentario ordinario. De hecho, el FMI y la OCDE han señalado que la proliferación de fondos extrapresupuestarios es uno de los indicadores más fiables de opacidad fiscal en cualquier país, precisamente porque permite eludir la disciplina que el principio de unidad pretende garantizar.

Conviene señalar que este principio no ha quedado libre de una corriente doctrinal de sentido contrario. Las reformas de Nueva Gestión Pública de los años ochenta y noventa (caso neozelandés) defendieron, precisamente, una mayor fragmentación deliberada del presupuesto en unidades ejecutoras más autónomas, dotadas de mayor flexibilidad de gestión a cambio de mayor responsabilidad sobre resultados medibles, bajo el argumento de que la rigidez de la unidad presupuestaria clásica podía resultar, paradójicamente, un obstáculo para la propia eficiencia que el principio correspondiente persigue. 

\paragraph{Especificación}
El principio de especificación impone una obligación dual sobre la actividad presupuestaria: un límite cuantitativo (no exceder créditos), y un límite cualitativo (no modificar su destino). 

Responde a la misma lógica que ya trabajamos con la clasificación funcional. Si el Parlamento aprueba una partida finalista, el Ejecutivo no puede, sin nueva autorización, redirigirla a otro fin. Es la garantía de que la aprobación parlamentaria no se convierte en un cheque en blanco disfrazado de autorización específica. El propio conflicto presupuestario prusiano que acabamos de narrar es, en realidad, tanto un episodio sobre el principio de competencia como sobre este: la disputa concreta entre Bismarck y el Landtag no era una discusión abstracta sobre quién aprueba el presupuesto, sino una disputa muy específica sobre el destino concreto de una partida de gasto.

\paragraph{No afectación}
Atendiendo al principio de no afectación, los ingresos forman una masa común, sin asignación a un gasto específico. JÈZE la formula como corolario directo del principio de universalidad: si los ingresos forman una masa común sin destino específico, es porque la propia universalidad presupuestaria exige que ningún ingreso pertenezca de antemano a ningún gasto concreto, preservando así la libertad de asignación del Parlamento en cada ejercicio.

Este principio tiene, detrás, un debate genuino de eficiencia asignativa. Si un ingreso concreto queda afectado (hipotecación o earmarking), se corre el riesgo de que el nivel de gasto en esa política quede determinado por la volatilidad de la recaudación de ese tributo concreto, en lugar de por una valoración autónoma de cuánto merece gastarse en esa política a la luz de todas las prioridades del sector público en su conjunto. 

Pero existe, frente a ese argumento, uno de sentido contrario y de peso político real. La afectación puede reforzar la legitimidad social de un tributo, al mostrar con claridad a quien lo paga adónde va su dinero, y es también el argumento que ha ido ganando peso en el diseño reciente de tributos ambientales, donde vincular explícitamente la recaudación a gastos de transición energética facilita su aceptación social frente a un impuesto verde percibido como puramente recaudatorio. Como ejemplo, este es el argumento que sostiene la afectación de las cotizaciones sociales a las pensiones contributivas, reforzando la percepción de cuota de un seguro frente a la de impuesto genérico. 

\paragraph{No compensación}
Conforme al principio de no compensación, ingresos y gastos deben figurar por su importe íntegro, sin compensarse entre sí. Igual que el anterior, JÈZE lo formula como manifestación del principio de universalidad presupuestaria: ingresos y gastos deben figurar por su importe íntegro y sin compensarse entre sí (\textit{règle de non-contraction}), precisamente porque permitir esa compensación equivaldría, de facto, a sacar del presupuesto la parte neteada. Exactamente el mismo riesgo de opacidad que el principio de unidad pretende conjurar frente a los fondos extrapresupuestarios.

No obstante, existen notables excepciones. Las devoluciones de impuestos se registran como menores ingresos, no como gasto. La razón de fondo es de transparencia. Si se permitiera netear ingresos y gastos relacionados, el volumen bruto de la actividad del sector público quedaría sistemáticamente infravalorado en la presentación presupuestaria, dificultando cualquier comparación seria de la verdadera magnitud del Estado dentro de la economía.

\subsection{Marco institucional de la actividad presupuestaria del Estado}
¿Quién, dentro del sector público español, ejerce esta actividad presupuestaria?

Conviene distinguir con precisión dos preguntas. La primera se pregunta por qué, dentro de un mismo nivel de gobierno, existen múltiples entidades jurídicas distintas del propio Estado-persona jurídica, cada una con su propio presupuesto: es una pregunta sobre técnica de organización administrativa. La segunda se pregunta por qué existen, además, distintos niveles territoriales de gobierno con potestad presupuestaria propia: es una pregunta sobre distribución territorial del poder. 

\subsubsection{Desagregación horizontal: Administración General y Administración Institucional}
Si el Estado ya es, de por sí, una persona jurídica capaz de actuar, contratar y presupuestar, ¿por qué necesita crear, además, decenas de personas jurídicas distintas de sí mismo para ejercer parte de su propia actividad? 

La respuesta que ofrece la doctrina clásica del Derecho administrativo español es la de la personificación instrumental:\footnote{%
    GARCÍA DE ENTERRÍA, en su distinción entre Administración General y Administración Institucional.
} crear una persona jurídica nueva, dotada de patrimonio y presupuesto propios pero dependiente en última instancia del Estado matriz, permite dotar a una actividad concreta de un régimen de gestión distinto del ordinario. Esta técnica, la descentralización funcional, responde a razones muy diversas según el tipo de entidad, y merece la pena distinguirlas una por una, porque no todas persiguen lo mismo.

El catálogo actual de estas entidades está regulado por la Ley 40/2015 (LRJSP). La LRJSP introdujo, por primera vez con esta denominación positivizada, el concepto de Sector Público Institucional, integrado por: (i) los organismos públicos vinculados o dependientes de la AGE, clasificados a su vez en organismos autónomos, entidades públicas empresariales y agencias estatales; (ii) las autoridades administrativas independientes; (iii) las sociedades mercantiles estatales; (iv) los consorcios; (v) las fundaciones del sector público; (vi) los fondos sin personalidad jurídica; y (vii) las universidades públicas no transferidas.\footnote{%
    Sociedades mercantiles estatales: naturaleza jurídico-privada plena, participación estatal mayoritaria, presupuesto estimativo (SEPI, ICO);, consorcios, entes de cooperación entre Administraciones distintas, o entre Administración y entidades privadas sin ánimo de lucro; fundaciones del sector público (aportación mayoritaria estatal); y fondos sin personalidad jurídica (FIEX, FONPRODE, FIEM).
} 

Como vemos, es una lista tasada y cerrada, precisamente para evitar el fenómeno de proliferación incontrolada de figuras jurídicas atípicas que había caracterizado las décadas anteriores: solo en el proceso de racionalización se suprimieron $2.194$ entidades de distinta naturaleza dentro del sector público español.\footnote{%
    Conviene detenerse en una figura que la propia LRJSP eliminó, porque su historia es, en sí misma, un ejemplo revelador de la tensión entre eficiencia de gestión y disciplina fiscal. Las agencias estatales, creadas por la Ley 28/2006, pretendían introducir en España el modelo de agencia autónoma con gestión por objetivos y financiación plurianual. No obstante, la propia exposición de motivos justificó su supresión alegando que no habían cumplido las expectativas para las que se crearon. Pero la explicación más plausible, señalada por la doctrina que ha estudiado el episodio, apunta a una causa más concreta y más reveladora. Las propias medidas de control del gasto público derivadas de la crisis financiera de 2008-2012 y de la LOEP resultaban, en la práctica, incompatibles con el grado de autonomía financiera y de gestión que el modelo de agencia exigía para funcionar según su diseño original. Un instrumento diseñado para ganar eficiencia de gestión termina siendo desactivado por chocar frontalmente con el reforzamiento de la disciplina fiscal centralizada que la propia crisis impuso.
}

\paragraph{Organismos Autónomos: gestión especializada bajo control jerárquico pleno}
Los organismos autónomos son entidades de derecho público, con personalidad jurídica, tesorería y patrimonio propios, dependientes de un departamento ministerial y con presupuesto limitativo; es decir, no pueden sobrepasar el gasto autorizado. 

La razón de crear un organismo autónomo es, fundamentalmente, de \textbf{especialización funcional}. Agrupar bajo una gestión unificada y con contabilidad propia una actividad que, aun siendo de naturaleza puramente administrativa, tiene una entidad y una complejidad de gestión que justifica dotarla de cierta autonomía operativa, sin llegar a salir del control jerárquico y presupuestario más estricto propio de la Administración General.

\paragraph{Entidades Públicas Empresariales: cuando la actividad se acerca al mercado}
La diferencia de fondo con el organismo autónomo es de naturaleza económica de la actividad. Una entidad pública empresarial desarrolla actividades susceptibles de contraprestación económica, con una lógica más próxima al tráfico mercantil privado (ADIF, ENAIRE o Puertos del Estado).

Esto justifica dotarlas de un presupuesto estimativo y de un régimen jurídico-privado más flexible para su gestión ordinaria, aunque conserven naturaleza pública y control último del Estado.

\paragraph{Autoridades Administrativas Independientes: el problema constitucional de la independencia dentro de un Gobierno que dirige la Administración}
Esta es, con diferencia, la categoría que plantea la pregunta jurídico-conceptual más profunda de todo el catálogo. 

El art. 97 CE encomienda al Gobierno la dirección de la Administración, lo que en principio implica una cadena de responsabilidad política plena ante las Cortes. 

Las autoridades administrativas independientes (CNMC, la CNMV, el Banco de España, el Consejo de Seguridad Nuclear, y la AIReF) rompen deliberadamente esa cadena. Se les dota, por ley, de independencia funcional frente al Gobierno, precisamente para el tipo de función que desempeñan. La justificación doctrinal de esta ruptura es que ciertas funciones exigen una credibilidad técnica que resultaría estructuralmente debilitada si quien las ejerce puede ser cesado o instruido por el Gobierno de turno según su conveniencia política del momento. La AIReF, de hecho, es el ejemplo más nítido posible de esta lógica aplicada al tema. Carecería casi por completo de sentido que fuera el propio Ministerio de Hacienda quien evaluara, sin ningún control externo, si el propio Ministerio de Hacienda ha cumplido sus propios objetivos de déficit. 

Pero esta misma lógica genera una tensión constitucional que la doctrina no considera resuelta de forma definitiva. ¿Hasta qué punto es compatible, en términos de legitimidad democrática, que decisiones con enorme impacto económico y social queden en manos de órganos técnicos cuyos responsables no son elegidos ni cesables por el circuito de responsabilidad política ordinario? Es el mismo debate que ya rozamos al hablar de la credibilidad de las reglas fiscales frente a la discrecionalidad democrática. Cuanto más se blinda una decisión frente al ciclo político para ganar credibilidad técnica, más se aleja también del circuito de responsabilidad democrática directa que legitima cualquier decisión pública.

\subsubsection{Desagregación vertical: Reparto competencial en materia presupuestaria}
Como sabemos, el Título VIII de la Constitución configura a España como un Estado profundamente descentralizado territorialmente, reconociendo autonomía a municipios, provincias y Comunidades Autónomas para la gestión de sus propios intereses. Pero conviene entender, con precisión jurídica, por qué la autonomía política que reconoce la Constitución requiere una autonomía financiera propia. 

Una Comunidad Autónoma o un municipio con competencias plenas sobre sanidad, educación o servicios sociales, pero sin ningún margen propio de decisión sobre sus ingresos, no sería autónoma en ningún sentido sustantivo. Es la misma intuición que sostenía el principio medieval de \textit{no taxation without representation}: el poder político sin poder financiero propio es un poder político incompleto. 

El art. 156.1 CE reconoce, por ello, explícitamente la autonomía financiera de las Comunidades Autónomas para el desarrollo y ejecución de sus competencias. El art. 142 CE, en cambio, emplea para las Haciendas Locales un término deliberadamente más modesto: suficiencia financiera, no autonomía plena.\footnote{%
    La razón de esta asimetría constitucional no es arbitraria. Responde a que las Comunidades Autónomas, especialmente las llamadas históricas, heredaban una tradición institucional y fiscal propia que el constituyente decidió respetar y actualizar, mientras que los municipios y provincias españoles, con la parcial excepción histórica de determinados regímenes forales locales, carecían de una tradición comparable de autonomía tributaria propia con ese mismo peso histórico, lo que justificó constitucionalmente un diseño de garantía financiera en lugar de un diseño de autonomía tributaria plena y genuina.
}

\paragraph{LOFCA: Régimen común}
Para las Comunidades Autónomas de régimen común, la ley orgánica de desarrollo de las competencias financieras es la Ley Orgánica 8/1980 (LOFCA), la norma marco que ha regulado los principios generales del sistema de financiación.

Conviene señalar, como dato de actualidad relevante, que el sistema de financiación autonómica de régimen común no ha sido un diseño estático. Desde 1980, España ha atravesado sucesivas renegociaciones de su modelo financiación, cada una intentando corregir los desequilibrios que el modelo anterior había revelado en la práctica (suficiencia, equidad interterritorial, incentivos a la corresponsabilidad fiscal). El modelo actualmente vigente, aprobado por la Ley 22/2009, lleva en vigor más de quince años, pese a que la propia norma prevé mecanismos periódicos de evaluación y actualización, y pese a que el propio debate político español ha señalado repetidamente, en los últimos años, la necesidad de una reforma que el modelo actual sigue sin recibir. 


Frente a este régimen común, el País Vasco y Navarra se rigen por un sistema radicalmente distinto: el Concierto Económico vasco y el Convenio Económico navarro, sistemas bajo los cuales son las propias instituciones forales —las Diputaciones Forales vascas, la Comunidad Foral de Navarra— quienes recaudan, gestionan e inspeccionan la práctica totalidad de los tributos, incluidos el IRPF y el Impuesto sobre Sociedades, dentro de su propio territorio, aportando después al Estado una cuota o aportación en compensación por los servicios y competencias que el Estado sigue prestando en su territorio o en beneficio del conjunto de España (defensa, política exterior, entre otros). [Seguro] El fundamento constitucional de esta asimetría, deliberadamente distinta del régimen común, se encuentra en la disposición adicional primera de la Constitución, que "ampara y respeta los derechos históricos de los territorios forales", previendo su "actualización general" dentro del marco de la propia Constitución y de los Estatutos de Autonomía. [Probable] La justificación jurídica de fondo no es, por tanto, una simple concesión política de conveniencia, sino el reconocimiento constitucional expreso de una continuidad histórica preconstitucional —los fueros vascos y navarros, con su propia tradición fiscal diferenciada, anteceden a la propia Constitución de 1978, y esta optó explícitamente por incorporarlos, actualizados, en lugar de disolverlos en el régimen común—. [Suposición] Esta asimetría, precisamente por su origen histórico y no puramente funcional, ha generado un debate político recurrente sobre su compatibilidad con el principio de solidaridad interterritorial del art. 138 CE —los críticos del sistema foral señalan que estos territorios, generalmente entre los de mayor renta per cápita de España, contribuyen proporcionalmente menos a la solidaridad interterritorial que las Comunidades de régimen común, precisamente por gestionar y quedarse la práctica totalidad de su propia recaudación—, un debate que, a diferencia de las reglas fiscales europeas que ya trabajamos, no tiene visos de resolverse por vía de reforma técnica, precisamente por descansar sobre un fundamento de legitimidad histórica y no meramente de eficiencia recaudatoria.

\paragraph{TRLRHL: Entidades Locales}
Para el ámbito local, el desarrollo del principio constitucional de suficiencia financiera del art. 142 CE se encuentra hoy en el TRLRHL, que regula los recursos propios de municipios y provincias (IBI, IAE, IVTM), tasas, contribuciones especiales, junto con la participación en los tributos del Estado. 

Esta arquitectura vertical de tres niveles (Estado, Comunidades Autónomas, Entidades Locales) más la propia Seguridad Social como cuarto subsector, es exactamente la misma que sostiene toda la LOEP y todo el aparato de reglas fiscales. No es casualidad que la propia distinción constitucional entre autonomía financiera plena (CCAA) y mera suficiencia financiera garantizada (EELL) se refleje en el reparto desigual de los límites de deuda que ya conocimos ($44\%$ AACC, $13\%$ CCAA, apenas $3\%$ EELL): a menor autonomía financiera menor el margen de endeudamiento propio que el legislativo les reconoce.


\clearpage
\section{Política Presupuestaria: Ciclo Presupuestario}
\puntosclave{
    Las distintas fases a las que se someten los PGE son la elaboración (poder ejecutivo). la aprobación (poder legislativo). la ejecución (poder ejecutivo) y el control ( los tres poderes) .
    
    • Las sub-fases de la elaboración del Proyecto de PGE son: aprobación de los objetivos fiscales y del techo de gasto. elaboración de presupuestos por pa11e de las entidades del sector público estatal, formación del Anteproyecto de PGE y elevación al Consejo de M inistros. aprobación del Proyecto de PGE. y remisión a las Cortes y a la UE.
    
    • La fase de aprobación en Cortes se resume en: comparecencia de altos cargos. debate de las enmiendas a la totalidad en el Pleno. debate de las enmiendas parciales en la Comisión de Presupuestos. y aprobación en el Pleno. Si se aprueba en el Congreso. se remite al Senado, donde el proceso es similar. Si el Senado rechaza el Proyecto de PGE o lo enmienda. vuelve al Congreso (que puede levantar tanio el veto como las enmiendas del Senado).
    
    • En la fase de ejecución del presupuesto, el Gobierno implementa las actuaciones cubienas por los créditos de los PGE, debiendo cefiirse a la finalidad para la que los créditos fueron aprobados.
    
    • Por último, en la fase de control de la ejecución del presupuesto intervienen los tres poderes: el ejecutivó (a través de la lG AE), el legislativo ( a través del Tribunal de Cuentas) y el judicial (a través de los órganos judiciales).
}


Cerrado ya el marco jurídico-conceptual, entramos en el núcleo del tema: el ciclo completo del presupuesto.

La Ley General Presupuestaria define los presupuestos como la expresión cifrada, conjunta y sistemática de los derechos y obligaciones a liquidar durante el ejercicio por cada uno de los órganos y entidades que forman parte del sector público estatal (art. $32$).

El Proyecto de Ley de PGE se compone de dos partes (art. 37):\footnote{El articulado legal recoge normas para la gestión de los créditos.
} el articulado legal y los estados contables (cifras de ingresos y gastos de los sujetos que integran el sector público estatal). Además, el Proyecto de Ley PGE se tiene que acompañar de otra documentación complementaria, como la liquidación de los presupuestos del año anterior, la memoria de beneficios fiscales, las inversiones, etc.

Las distintas fases a las que se someten los PGE son la elaboración, la aprobación, la ejecución y el control. La elaboración y la ejecución recaen en el poder ejecutivo, la aprobación la lleva a cabo el poder legislativo, y el control lo ejercen los tres poderes ( ejecutivo. legislativo y judicial).

\subsection{Elaboración}
La responsabilidad principal de la elaboración de los Presupuestos recae sobre el Ministerio de Hacienda, pero requiere la participación del conjunto del Gobierno y de la Administración General del Estado. 

\subsubsection{Objetivos Fiscales: }


\paragraph{Unión Europea: Semestre Europeo}
El ciclo de elaboración de los PGE no arranca, como cabría esperar, con una decisión puramente doméstica del Gobierno español. El ciclo presupuestario nacional se encuadra en el ciclo europeo denominado Semestre Europeo. Este engloba el conjunto de actuaciones de las instituciones europeas y de los Estados miembros en relación con la gobernanza fiscal y económica. El Semestre Europeo es, por tanto, un ciclo de coordinación de las políticas económicas, presupuestarias, sociales y de empleo dentro de la UE.

Entonces, el ciclo de elaboración arranca con el Semestre Europeo, el ciclo anual de coordinación de las políticas económicas que se inicia en noviembre con el Estudio Prospectivo Anual sobre el Crecimiento Sostenible de la Comisión y continúa en febrero con los Informes-país.\footnote{%
    Conviene entender por qué existe este mecanismo. El Semestre Europeo se creó en 2010-2011, como respuesta directa al fracaso constatado del PEC original de 1997 para prevenir la crisis de deuda soberana europea. El PEC original operaba con una lógica esencialmente ex post. Los Estados miembros presentaban sus Programas de Estabilidad, y la Comisión los evaluaba y, en su caso, sancionaba, después de que las decisiones presupuestarias nacionales ya estuvieran tomadas. Este diseño se reveló insuficiente e incapaz de lograr en forma efectiva sus objetivos, precisamente porque llegaba tarde, cuando la Comisión formulaba su recomendación, el presupuesto nacional del año en cuestión ya estaba cerrado. El Semestre Europeo invierte esa lógica, desplazando la supervisión europea a un momento ex ante. Antes de que cada Gobierno cierre su propio Proyecto de Presupuestos, ya ha tenido que remitir su Programa de Estabilidad, recibir las recomendaciones específicas por país del Consejo, y remitir también su Plan Presupuestario en octubre para que la Comisión emita opinión sobre él antes de su aprobación parlamentaria definitiva. Este cambio de secuencia no fue una decisión aislada. Se enmarca en el mismo paquete legislativo de reforma de la gobernanza económica europea, el Six-Pack de diciembre de 2011, que es el mismo contexto normativo europeo que empujó a España a reformar su art. $135$ CE apenas unos meses antes, en septiembre de ese mismo 2011.
}

A finales de abril el Gobierno procede a la Actualización del Programa de Estabilidad (PE). Este documento recoge las previsiones macroeconómicas y fiscales para el año en curso y para los tres ejercicios posteriores. El PE constituye en España el documento de planificación presupuestaria a medio plazo al que hace referencia la Directiva de Marcos Fiscales Nacionales.\footnote{%
    La AIReF señala que el contenido del PE no es suficiente para suponer una herramienta útil y eficaz para la coordinación y planificación a medio plazo de la política presupuestaria, por dos motivos. Primera, porque el desglose de ingresos y gastos en términos de contabilidad nacional solo se ofrece para el conjunto de las AAPP, sin desagregar por subsectores, de modo que el documento que en teoría debería servir de marco de planificación conjunta del Estado autonómico no permite, en realidad, saber qué se espera de cada uno de sus niveles. Segunda, y más grave desde el punto de vista de la lealtad institucional, el Programa de Estabilidad lo elabora el Gobierno central con la información que le remiten las Administraciones Territoriales, sin que estas participen ni sean corresponsables del documento final que, sin embargo, condicionará después sus propios objetivos fiscales individuales. 

En cierto sentido, es la misma tensión que ya identificamos con la restricción presupuestaria blanda. Un sistema en el que el nivel central diseña unilateralmente el marco de disciplina fiscal que después se impone, mediante el CPFF, a los niveles territoriales, plantea un problema genuino de legitimidad democrática de la coordinación vertical, incluso cuando el resultado técnico final sea razonable en su contenido.
}

El PE es remitido a las autoridades europeas antes del 30 de abril de cada año. Antes suele discutirse en la Comisión Delegada de Asuntos Económicos (CDGAE) del Gobierno. Además, la AIReF debe emitir antes del 15 de abril un informe sobre el PE. Sin embargo, en la práctica el Gobierno agota el plazo y suele enviar el PE a Europa el mismo 30 de abril, por lo que la AIReF realiza el informe posteriormente. La razón de esto es que hasta finales de marzo no se conoce el avance del déficit del año anterior, por lo que este incumplimiento formal queda convalidado con la aprobación posterior de los objetivos fiscales. Lo que sí tiene que emitir la AIReF con carácter previo a la elaboración del PE es el aval sobre el cuadro macroeconómico de referencia elaborado por el Gobierno. 

Tras ello, la Comisión analiza el PE y publica en mayo una evaluación del mismo, así como una propuesta de recomendaciones específicas por país. En juni o julio, el Consejo de la Unión (ECOFIN), aprueba las recomendaciones específicas por país propuestas por la Comisión.

\paragraph{Ejecutivo Español}
Tras la remisión del Programa de Estabilidad y la valoración por parte de la Comisión Europea, el Gobierno inicia el procedimiento de los artículos $15$ y $16$ de la LOEP para aprobar los objetivos fiscales. Sin embargo, es costumbre avanzar previamente a la ratificación del ECOFIN, ya que, en la práctica, suelen aprobarse sin modificaciones sustanciales.

Los objetivos de déficit y de deuda y la regla de gasto se aprueban para los tres años siguientes y para cada uno de los cuatro subsectores que componen las Administraciones Públicas. Paralelamente a la elaboración de éstos, el Gobierno elabora también el límite de gasto no financiero del Estado para el año siguiente (conocido como techo de gasto).\footnote{%
    Este paso es necesario porque, aunque los objetivos fiscales están denominados en contabilidad nacional (devengo puro), los presupuestos se definen en contabilidad pública o presupuestaria (devengo imperfecto, más próximo a la caja), por lo que para trasladar los objetivos fiscales del Estado a los PGE es necesario realizar algunos ajustes, para lo cual se elabora el techo de gasto.
}

Antes de la aprobación por el Consejo de Ministros, el Gobierno (Ministerio de Hacienda) tiene que informar al CPFF sobre los objetivos fiscales de su subsector (antes del 1 de abril) y del techo de gasto del Estado (antes del 1 de agosto). El Gobierno suele cumplir con este mandato legal con una única reunión del CPFF en verano. Como para la elaboración del techo de gasto el Gobierno ya ha cuantificado el sistema de financiación autonómico para el año próximo, el Ministerio de Hacienda aprovecha también la reunión para comunicar a las CCAA los recursos de los que van a disponer por las dos vías del sistema de financiación (impuestos cedidos, participación y fondos), de manera que con esa información las CCAA ya pueden comenzar a elaborar sus presupuestos autonómicos para el ejercicio siguiente.

En la reunión, el Gobierno también facilita el conocido como Informe de Situación, un informe elaborado por el Ministerio de Economía que contiene el escenario macroeconómico, la tasa de referencia para la regla de gasto, la posición cíclica de la economía y el saldo cíclico de las Administraciones Públicas en su conjunto y por subsectores, entre otras cosas. Por último, el Ministerio de Hacienda se reúne también con la Comisión Nacional de la Administración Local (CAL) para informar de los objetivos fiscales de su subsector y de los recursos del sistema de financiación local.

Tras el innforme del CPFF y de la CNAL, el Consejo de Ministros aprueba los objetivos fiscales para cada una de las Administraciones. En ese mismo Acuerdo del Consejo de Ministros se aprueba el techo de gasto antes mencionado.

\paragraph{Cortes Españolas}
Una vez aprobados por el Gobierno, los objetivos de déficit y de deuda (no la regla de gasto) deben ser ratificados o rechazados por el Congreso y el Senado.\footnote{%
    Nótese que las Cortes aprueban exclusivamente los objetivos de déficit y de deuda, no la regla de gasto (automática) ni el techo de gasto del Estado, quizá por no ser necesario en la medida en que este se deriva mecánicamente del objetivo de déficit. No obstante, los acuerdos del Congreso y del Senado aprobando los objetivos de déficit y deuda también suelen mencionar el techo de gasto.
} En la documentación que acompañe al acuerdo de fijación de objetivos también se incluyen las recomendaciones específicas a España de la Unión (CSRs) y el Informe de Situación.

Si alguna de las dos cámaras rechaza los objetivos de déficit y deuda, el Gobierno debe reiniciar el proceso. Al contrario de lo que ocurre en la tramitación legislativa ordinaria, no hay segunda lectura de los objetivos en el Congreso. 

Por ejemplo, los objetivos fiscales de 2019-2018, la primera propuesta remitida fue rechazada por el Congreso y la segunda fue rechazada por el Senado. Tras esto, el Gobierno decidió no enviar de nuevo a las Cortes los objetivos fiscales y optó por elaborar los PGE de 2019 con los objetivos antiguos para 2019 que se habían aprobado en 2017, PGE 2018. Conviene situar este episodio dentro de un fenómeno más amplio y de mayor actualidad. La fragmentación parlamentaria que ha caracterizado la política española desde 2015 ha convertido la aprobación de nuevos PGE en un obstáculo recurrente. España ha atravesado, en los últimos años, periodos prolongados de prórroga presupuestaria precisamente por la incapacidad de sucesivos Gobiernos en minoría de reunir la mayoría parlamentaria necesaria para aprobar un Proyecto de Ley de Presupuestos nuevo. Este fenómeno, que la Constitución previó como una válvula de seguridad excepcional, se ha convertido en un régimen casi ordinario de gestión presupuestaria durante buena parte de una legislatura entera. Una transformación que ninguno de los constituyentes de 1978, ni tampoco los reformadores de 2011, previeron como escenario habitual, y que plantea una pregunta que la propia norma no resuelve del todo: ¿hasta qué punto sigue siendo compatible con el principio de anualidad un sistema en el que ese consentimiento no se renueva?

\paragraph{Consejo de Política Financiera y Fiscal}
Tras la aprobación de los objetivos globales por el Consejo de Ministros, el Ministerio de Hacienda, previo Informe de la AIReF, realiza una propuesta individualizada de objetivos de déficit y de deuda para cada una de las Comunidades Autónomas. Con esta propuesta, y previo informe del Consejo de Política Fiscal y Financiera, el Consejo de Ministro aprueba los objetivos de déficit y de deuda individualizados para cada CCAA. En la práctica, sin embargo, no se suelen formular objetivos de déficit diferenciados para cada CCAA.

Conviene señalar que el Ministro de Hacienda ostenta el $50\%$ de los derechos de voto del CPFF en su conjunto, frente al otro $50\%$ repartido entre las diecisiete Comunidades Autónomas. Un reparto de voto que, en la práctica, otorga al Gobierno central una posición de veto efectivo sobre cualquier decisión que el CPFF pudiera adoptar en contra de su criterio, con independencia de cuántas Comunidades Autónomas se opongan.\footnote{%
    Es un dato que las propias Comunidades Autónomas han señalado, en distintos momentos y con distinto signo político según qué partido gobernara en Madrid, como una de las asimetrías más visibles de todo el sistema de coordinación fiscal territorial español.
} Por tanto, se trata de un órgano formalmente multilateral, cuya arquitectura de voto garantiza que el nivel central nunca pueda quedar en minoría.

\subsubsection{Presupuestos Generales del Estado}
La elaboración del Proyecto de PGE se inicia con la publicación de la Orden Ministerial por la que se dictan las normas para la elaboración de los PGE.

Esta Orden recoge los criterios generales de presupuestación, el ámbito institucional, el proceso de elaboración y tramitación, la definición de los órganos participantes, los plazos y documentación, y la descripción de las estructuras presupuestarías. Esta orden es complementada a su vez por diversas Resoluciones de la Dirección General de Presupuestos, departamento responsable de la elaboración del Anteproyecto de Ley.

La Comisión de Políticas de Gasto es presidida por el Ministro de Hacienda y está integrada por todos los Ministros. La Comisión marca las líneas principales de distribución del gasto y la asignación de cantidades por políticas (sanidad. educación, agricultura. digitalización. vivienda. etc.).


Las distintas unidades del sector público estatal elaboran sus anteproyectos parciales de presupuestos, que después remiten a la Dirección General de Presupuestos.

La técnica empleada por la legislación española desde los af'íos 80 es la presupuestación del gasto por programas, donde la unidad principal de análisis y evaluación ·son los programas de gasto, que se definen como un conjunto de créditos presupuestarios destinados a un mismo fin (p.ej.  mantenimiento de caneteras). La presupuestación por programas tiene un enfoque por resu ltados. por lo que requiere de la definición de objetivos y de indicadores para la evaluación. De ahí que las unidades del sector público estatal. a la hora de elaborar sus presupuestos, no tengan que definir solamente las cifras de ingresos y gastos, sino también el aborar memorias ele programas donde describen el programa, sus objetivos y los indicadores a emplear.\footnote{%
    No obstante, en la práctica. la elaboración (y ej ecución) presupuestaria suele tener un enfoque más orgánico ( qué Ministerio) que funcional (qué programas). conservando aún ci erto carácter incrementalista y con una definición de objetivos e indicadores con margen de mejora, de manera que el seguimiento del gasto no se realiza tanto desde la óptica de la consecución de los objetivos e indicadores. sino del grado de ej ecución presupuestaria.
}

Por su parte, los ingresos del Estado los elabora la Secretaría de Estado de Hacienda, y dependen del ciclo económico. del marco tributario vigente, y de las medidas discrecionales que se contengan en los PGE.

La Dirección General de Presupuestos es el centro responsable de la elaboración del Anteproyecto de los PGE. Para ello. se coordina con las Oficinas Presupuestarias que se encuadran dentro de la Subsecretaría de cada Ministerio. Estos centros se encargan de la elaboración y gestión de los presupuestos de su Ministerio y canalizan la información hacia la Dirección General de Presupuestos. que agrupa los anteproyectos de las distintas unidades.

Una vez recibidos los anteproyectos parciales de presupuestos de los Ministerios, normalmente en verano, se produce el proceso de negociación interna entre los Ministerios sectoriales y el Ministerio de Hacienda.

Para canalizar este debate y fomentar la asignación eficiente de los recursos, por cada Ministerio se constituye una Comisión de Análisis de Programas. Son presididas por el Secretario de Estado de Presupuestos y Gastos (o, en su ausencia, por el Director General de Presupuestos) e integradas por el Subsecretario del Ministerio correspondiente y por el jefe de su Oficina Presupuestaria. así como de otros responsables de gasto que se consideren rele vantes. Tienen como fu nción analizar los anteproyectos de cada Ministerio desde un enfoque de programas, así como su encaje dentro de las restricciones presupuestarias definidas previamente.



Los anteproyectos de los diferentes Ministerios se van ajustando en este proceso iterativo que no concluye de manera definitiva hasta que el Ministerio de Hacienda aglutina los anteproyectos de los centros conforma el Anteproyecto de Ley PGE y lo eleva al Consejo ele Ministros para su aprobación (normalmente a la vuelta del verano.\footnote{%
    Los presupuestos de la Seguridad Social son elaborados por el Ministerio de Seguridad Social. Estos presupuestos incluyen los del Instituto de Gestión Sanitaria (INGESA), los del Instituto de Mayores y Servicios Sociales (lMSERSO). y los de las Entidades Gestoras, Servicio Común y Mutuas de Accidentes de Trabajo y Enfermedades Profesionales. El Ministerio de Seguridad Social remite su Anteproyecto al Ministerio de Hacienda para su integración con los PGE, y ambos ministros elevan conj untamente el Anteproyecto de Presupuesto de la Seguridad Social al Consejo de Ministros para su aprobación.
}

Conviene tener en cuenta que el margen de discrecionalidad real del Gobierno al diseñar este Proyecto es muy reducido: menos de $10.000$ millones de euros de gasto genuinamente nuevo, una vez descontados los compromisos ineludibles. Evidentemente, relativiza la propia intensidad del debate político que cada año rodea la negociación de los PGE, pues la inmensa mayoría de lo que se está aprobando cada año no es objeto de decisión nueva, sino la continuidad casi automática de compromisos ya asumidos en ejercicios anteriores. También hay que señalar que existen disposiciones de los PGE que tienen incidencia más allá del año en el que son aprobadas. Así, por ejemplo, cuando se regula la tasa de reposición del empleo público o el incremento de los salarios públicos y pensiones públicas, se está condicionando tanto el gasto del ejercicio corriente como su evolución a medio y largo plazo.

Tras el verano, el Consejo de Ministros aprueba el Proyecto de Ley de Presupuestos Generales del Estado para el año siguiente.

Una vez aprobado, el Gobierno debe remitir a las Cortes el Proyecto de Ley PGE y la documentación complementaria antes del I de octubre. No obstante, por diversos motivos de índole política y económica, este plazo se ha incumplido desde 20 1723.

El día de la remisión24 , el Gobierno tambi��n publica el libro de presentación del Proyecto de PGE (también conocido como 'libro amarillo '), en el que se explica el Proyecto de PGE con un tono divulg ativo25. La remisión del Proyecto de ley PGE se for maliza con la entrega de la documentación por parte del Ministro de Hacienda al Presidente del Congreso y la posterior rueda de prensa del Ministro y publicación de toda la documentación del Proyecto. El libro amarillo es la forma más senci lla de aproximarse a los PGE.


\begin{specialblock}{\textbf{Documentación}: ¿qué documentación se remite a Cortes?}
    Toda la documentación del Proyecto que el Gobierno remite a las Cortes se organiza en cinco paquetes:
    \begin{lnum}
        \item \textbf{Serie roja}. Incluye el articulado de la Ley y las cifras de ingresos y gastos (es decir, el Proyecto de Ley PGE propiamente dicho). En esta serie se contiene la clasificación funcional: la que explica en qué se gasta (programas de gasto). Cada uno de estos programas de gasto viene acompañado de una memoria descriptiva de los objetivos que persigue, así como de unos indicadores de seguimiento y evaluación;
        \item \textbf{Serie verde}. Incluye las clasificaciones orgánica (quién) y económica (cómo) de los gastos e ingresos. También incluye las cifras de inversión pública (territorializada a nivel de provincia) y los anexos de personal (que muestran el número de empleados públicos de las entidades del sector público con presupuesto limitativo;
        \item \textbf{Serie amarilla}. Incluye el informe económico y financiero, que explica las cifras del Proyecto. Estos se solapan con el libro amarillo. Esta serie también incluye la memoria de beneficios fiscales, las cifras del presupuesto consolidado (que son las relevantes a efectos del debate económico y político) y un anexo con los flujos financieros entre España y la Unión Europea;
        \item \textbf{Serie gris}. Incluye la liquidación de los PGE del año anterior y el avance de liquidación de los PGE del año en curso;
        \item \textbf{Serie azul}. Incluye cuatro informes transversales sobre los PGE: género, infancia, adolescencia y familia, alineamiento con los ODS y alineamiento con la transición ecológica.
    \end{lnum}
\end{specialblock}

La remisión a Bruselas, y el cierre del círculo con el que empezamos

El ciclo de elaboración se cierra, antes del 15 de octubre, con la remisión a la Comisión Europea y al Eurogrupo del Plan Presupuestario (Draft Budgetary Plan). Por tanto, no es el presupuesto aprobado el que Europa evalúa, sino el proyecto, antes incluso de que las Cortes españolas hayan concluido su propia tramitación parlamentaria. 

En el marco del Semestre Europeo, los Gobiernos de los Estados de la Eurozona deben remitir a la Comisión Europea y al Eurogrupo, antes del 15 de octubre, el Plan Presupuestario (Druft Budgetary Plan, DBP), que debe contener el Proyecto de PGE y los proyectos de presupuestos del resto d e Administraciones Territoriales (CCAA y EELL). así como un detalle de las medidas discrecionales contenidas en ellos.

Tras ello, y antes del 30 de noviembre, la Comisión publica su opinión sobre el Plan Presupuestario de cada país, evaluando el cumplimiento de las reglas fiscales. Por su parte. la AlReF emite el "Informe sobre los Proyectos y Líneas Fundamentales de Presupuestos de las AAPP". en el que se pronuncia sobre el cumplimiento de las reglas fiscales para el afio siguiente.


\subsection{Aprobación}
Tras la remisión a las Cortes por parte del Gobierno, se inicia su tramitación parlamentaria. Goza de una tramitación parlamentaria singular respecto de cualquier otra ley. En particular, tiene \textbf{prioridad} en su tramitación frente al resto de iniciativas legislativas, y sigue un \textbf{procedimiento} con plazos y reglas \textbf{propias} que el Reglamento del Congreso reserva en exclusiva para él.\footnote{%
    La justificación de esta prioridad no es muy distinta a la que vimos en la tramitación de la reforma constitucional (art. 135 CE) en 2011. El Estado no puede funcionar sin presupuesto aprobado y cuanto más se demore, mayor es el riesgo de tener que recurrir a la prórroga. Pero esta misma prioridad plantea una tensión con el principio de deliberación parlamentaria, pues cuanto más se acelera y se prioriza un debate, menos espacio queda para el escrutinio detallado que en teoría justifica la propia existencia del control parlamentario del gasto.
}

El proceso se inicia con la comparecencia de los altos cargos ante la Comisión de Presupuestos del Congreso. Asimismo, también suelen comparecer para dar una visión más general del contexto económico y de los propios PGE el Gobernador del Banco de España y el Presidente de la AIReF. Tras ello, los grupos parlamentarios presentan enmiendas, tanto a la totalidad como parciales. 

Las enmiendas a la totalidad reflejan un rechazo total del Proyecto y son debatidas y votadas en el Pleno del Congreso. Si se aceptan, suponen la devolución del Proyecto al Gobierno. Si son todas rechazadas, el Proyecto se remite a la Comisión de Presupuestos para continuar la tramitación. 

Dentro del mismo plazo también se presentan las enmiendas parciales, aunque existen limitaciones para su admisión a trámite. Cualquier incremento de gasto que se proponga deberá ir acompañado por una disminución en la misma sección presupuestaria (mismo Ministerio). El Gobierno podrá rechazar las enmiendas que supongan aumento de gastos o disminución de ingresos; y no sólo en el ámbito de la tramitación de los PGE, sino también con motivo de proposiciones de ley de los grupos parlamentarios. Las enmiendas parciales son debatidas en primer lugar en el seno de la Comisión de Presupuestos. Para ello, se constituye un grupo reducido denominado Ponencia, que aprueba las enmiendas sobre las que hay consenso suficiente y se reflejan en el Informe de la Ponencia. Posterionnente, las que no han sido analizadas se debaten y votan en el Pleno de la Comisión. 

\begin{specialblock}{\textbf{Enmiendas:} ¿por qué tienen un carácter tan peculiar?}
    Hemos señalado dos límites concretos a las enmiendas parciales: cualquier incremento de gasto debe compensarse con una disminución en la misma sección presupuestaria, y el Gobierno puede rechazar las enmiendas que supongan aumento de gasto o disminución de ingresos, no solo en la tramitación de los propios PGE, sino en cualquier proposición de ley posterior.

    Esta facultad, que a primera vista puede parecer una anomalía del sistema español, es en realidad una traducción literal de la tradición comparado europeo. El art. $40$ de la Constitución francesa de 1958 (V República), declara inadmisibles (ni siquiera debatibles) las proposiciones y enmiendas parlamentarias que tengan como consecuencia crear o agravar una carga pública, o disminuir los recursos públicos. De hecho, no se permite siquiera la compensación entre aumento y disminución de gastos dentro de la propia iniciativa parlamentaria, práctica que sí toleraban las Repúblicas anteriores, y la prohibición se extiende a la totalidad de los textos legislativos, no solo a la Ley de Presupuestos.
    
    La razón histórica de esta institución, siguiendo a CARCASSONNE, es clara y se formula sin ambages. La Constitución de 1958 puso, en materia financiera, al Parlamento francés bajo tutela, como reacción directa a la experiencia de las Repúblicas Tercera y Cuarta, durante las cuales los diputados tenían tendencia sistemática a agravar los presupuestos mediante enmiendas que aumentaban el gasto sin ninguna contrapartida de responsabilidad política sobre el déficit resultante. El mismo problema, que la literatura económica reciente (ALESINA y PEROTTI, 1996) documentaría después con rigor empírico bajo el nombre de problema del pozo común (\textit{common pool problem}): cuando múltiples actores pueden proponer gasto adicional sin asumir individualmente el coste fiscal completo de esa decisión, el resultado sistemático es un sesgo estructural hacia el sobregasto, mismo fenómeno de indisciplina fiscal descentralizada que trató RODDEN y la restricción presupuestaria.\footnote{%
        El coste total de dichas peticiones se repartiría entre todos los contribuyentes futuros, y ningún grupo parlamentario concreto carga en solitario con la responsabilidad del déficit agregado resultante.
    }
    
    Por tanto, vemos como España no importó esta institución por casualidad. El constituyente, adoptó de forma explícita el modelo de racionalización que Francia ya había ensayado veinte años antes, precisamente para evitar que un Parlamento fragmentado pudiera desestabilizar las cuentas públicas mediante el ejercicio irresponsable del derecho de enmienda. El art. 134.6 CE es, así, una pieza más de ese mismo empeño constituyente de 1978 por blindar la estabilidad institucional frente a los riesgos que la propia historia española —y la francesa, ya vivida y superada veinte años antes— había demostrado sobradamente reales.

    Ahora bien, hay que señalar el reverso de esta lógica. Si el Parlamento no puede, mediante enmienda, aumentar el gasto total ni disminuir los ingresos previstos, y solo puede redistribuir dentro de un sobre presupuestario ya fijado por el Gobierno, entonces el margen real de control parlamentario sustantivo sobre los PGE queda muy acotado. 
    
    Las Cortes pueden decidir cómo se reparte el gasto dentro del total que el Gobierno ha decidido, pero no pueden decidir de forma autónoma cuánto debe gastar el Estado en su conjunto ni de dónde deben proceder sus ingresos. Es la misma paradoja que señalamos con la ley ómnibus de acompañamiento, una institución diseñada para proteger la disciplina fiscal frente al riesgo de indisciplina parlamentaria termina reduciendo el papel del Parlamento en el proceso presupuestario a algo más cercano a la ratificación cualificada que al control sustantivo pleno que el propio principio de anualidad parece prometer.
\end{specialblock}

El texto resultante tras incorporar las enmiendas aprobadas se recoge en el Dictamen de la Comisión y se eleva al Pleno. En el Pleno del Congreso se debaten aquellas enmiendas que los grupos parlamentarios deciden defender. 

Tras la aprobación en el Congreso, el texto se remite al Senado, donde la tramitación es casi idéntica, con las llamadas enmiendas de veto en lugar de enmiendas a la totalidad. Si el Senado veta el texto o lo enmienda, el Congreso puede levantar el veto por mayoría absoluta en primera votación o por mayoría simple en segunda votación dos meses después, y también puede ratificar las enmiendas del Senado o rechazarlas. Este diseño, que otorga al Congreso la última palabra, no es una peculiaridad de la tramitación presupuestaria, sino uno de los ejemplos más claros del debate aún abierto sobre la la naturaleza imperfecta del bicameralismo español.\footnote{%
    La doctrina ha señalado repetidamente que el Senado español carece en la práctica de un peso institucional equivalente al de las segundas cámaras de otros Estados descentralizados o federales comparados, precisamente porque en la inmensa mayoría de los procedimientos legislativos, incluido el presupuestario, el Congreso conserva la capacidad última de imponer su criterio. La tramitación presupuestaria es un caso más de un problema de diseño institucional que sigue sin resolverse. Un Senado que debería ser la cámara natural de coordinación territorial de la política fiscal (función que el CPFF desempeña fuera del circuito parlamentario), pero que en la práctica interviene en el proceso presupuestario con una capacidad de veto final que el propio diseño constitucional le niega.
}

Como sabemos, si a 1 de enero no hay PGE aprobados, se prorrogan automáticamente los créditos iniciales del ejercicio anterior, sin las modificaciones presupuestarias que hubieran podido sufrir a lo largo de aquel ejercicio, y sin los créditos destinados a gastos que terminen en el ejercicio ya cerrado (una obra pública concreta, por ejemplo). 

Esta doble exclusión no es casual. La prórroga es, por diseño, un mecanismo de continuidad mínima del Estado, no un instrumento para perpetuar decisiones de gasto puntuales o ya agotadas y tampoco para prorrogar modificaciones presupuestarias que, por su propia naturaleza excepcional, carecen de la vocación de continuidad que sí tienen los créditos iniciales ordinarios.\footnote{%
    Por ejemplo, sería estúpido prorrogar el crédito de una obra ya finalizada, porque ese gasto ya cumplió su función (créditos a gastos). Pero también sería contrario a la finalidad de la actividad presupuestaria
} 

Pero esta válvula de seguridad, se ha convertido en un régimen de gestión sostenido durante periodos mucho más largos de lo que el propio diseño constitucional parece haber anticipado. Una circunstancia que, por chocar con el espíritu del principio de anualidad, sigue generando un debate doctrinal no resuelto sobre si el marco jurídico actual necesita, o no, mecanismos adicionales que incentiven, u incluso obliguen, a una renovación más frecuente del consentimiento presupuestario parlamentario, en lugar de permitir que la prórroga se convierta en la opción por defecto.

\begin{specialblock}{\textbf{OJO ->} Aquí sería interesante introducir el debate sobre el papel del legislativo en la actividad presupuestaria}
    
    Todo sistema constitucional que exige aprobación periódica del presupuesto se enfrenta a la misma pregunta: ¿qué ocurre si, llegada la fecha límite, ese presupuesto no se ha aprobado? España, como la inmensa mayoría de los sistemas parlamentarios europeos, responde con la prórroga automática. Estados Unidos responde, en cambio, con el cierre administrativo: agotado el plazo sin apropiación aprobada, la Administración federal está legalmente obligada a cesar sus operaciones no esenciales, con empleados públicos enviados a suspensión de empleo sin sueldo hasta que el Congreso apruebe nueva financiación. Son, literalmente, las dos respuestas polares posibles, y conviene entender que no es casual: cada una es la traducción presupuestaria exacta de una arquitectura constitucional distinta.
    
    El origen jurídico del shutdown: una ley de 1884, reinterpretada en 1980
    
    Para empezar, conviene desmontar la idea de que el shutdown estadounidense es una consecuencia natural y siempre presente del sistema constitucional americano. No lo es: es el resultado de una reinterpretación jurídica muy concreta y muy tardía. 
    
    La base legal es la Antideficiency Act, una norma que prohíbe a las agencias federales gastar u obligar fondos sin apropiación del Congreso, originada en 1870 y codificada en su forma moderna en 1884; precisamente en el mismo contexto de disciplina fiscal de posguerra civil que ya vimos dar origen a instituciones de control del gasto. Pero durante prácticamente un siglo, esta ley se aplicó con una tolerancia notable. Entre 1950 y 1980 se produjeron al menos siete periodos sin apropiación aprobada, y en ninguna de ellas el Gobierno federal llegó a cerrar sus puertas. Las agencias simplemente continuaban operando, bajo el supuesto tácito de que el Congreso no pretendía realmente interrumpir el funcionamiento del Estado por una disputa presupuestaria puntual. De hecho, en varios de esos episodios, el Congreso ratificó a posteriori el trabajo realizado por empleados públicos durante el vacío, y se les pagó con normalidad.
    
    Todo cambió en 1980. Ante una laguna de financiación que afectaba a la Comisión Federal de Comercio (FTC) (producto de una disputa sobre un proyecto de ley que pretendía limitar sus competencias), el presidente Carter solicitó al Fiscal General Civiletti un dictamen formal sobre cómo debía interpretarse la Antideficiency Act. Civiletti concluyó que la ley era inequívoca. Durante los periodos de apropiaciones caducadas, no puede gastarse ningún fondo salvo el necesario para llevar a cabo la terminación ordenada de las funciones de la agencia, y que cualquier otro gasto constituiría una violación de la propia Antideficiency Act. Ese mismo día se produjo el primer cierre administrativo de la historia de Estados Unidos. Un solo día, 1.600 empleados suspendidos, un coste de 700.000 dólares, la mayor parte en horas de trabajo perdidas. Un segundo dictamen de Civiletti, en enero de 1981, precisó la llamada excepción de vida y propiedad: solo pueden continuar trabajando, sin cobrar hasta la resolución del vacío, los empleados cuya función guarde una conexión razonable y articulable con la seguridad de la vida humana o la protección de la propiedad; criterio que hoy sigue determinando qué funcionarios federales son esenciales y cuáles no. Por tanto, un mecanismo percibido como estructural es, en realidad, producto de la interpretación de un Fiscal General. Durante un siglo, el mismo marco jurídico había convivido con una práctica completamente distinta.
    
    Por qué esta diferencia no es casual: separación de poderes frente a fusión parlamentaria
    
    La razón de fondo por la que España, y la práctica totalidad de los sistemas parlamentarios europeos, nunca ha desarrollado nada parecido a un shutdown, no es una cuestión de tradición administrativa aleatoria, sino una consecuencia directa de la arquitectura constitucional de cada sistema. En un sistema parlamentario como el español, el Gobierno existe porque cuenta con la confianza del Parlamento, y esa confianza se expresa en la propia aprobación de su presupuesto. Si el desacuerdo entre Gobierno y Cortes sobre el presupuesto llegara a ser tan profundo como para bloquear indefinidamente su aprobación, el sistema dispone de otros mecanismos constitucionales para resolver esa crisis de confianza, en lugar de dejar que el propio funcionamiento cotidiano del Estado quede rehén de esa disputa. La prórroga es una solución de continuidad administrativa que presupone que la crisis política se resolverá por otras vías propiamente políticas.
    
    En un sistema presidencial como el estadounidense, Legislativo y Ejecutivo se eligen por separado, pueden pertenecer a partidos distintos, y ninguno depende de la confianza del otro para seguir existiendo. El presidente no puede ser destituido por el Congreso salvo mediante el procedimiento extraordinario de impeachment (pero sería de carácter personal, por lo que el VP sería el nuevo presidente), y el Congreso no puede ser disuelto por el presidente. En ausencia de esas válvulas de escape parlamentarias, el propio poder de la bolsa se convierte en el instrumento de choque genuino entre ambos poderes: si el Congreso no aprueba la apropiación que el presidente necesita para que la Administración funcione, no hay ningún mecanismo alternativo que resuelva automáticamente esa disputa, porque el diseño constitucional americano no quiso dotar a ninguno de los dos poderes de un instrumento para forzar al otro a ceder por otra vía que no fuera la negociación bajo la presión del propio cierre.
    
    La vía intermedia menos conocida: por qué la mayoría de los "vacíos" de financiación en EEUU no terminan en cierre
    
    Conviene matizar la imagen del sistema estadounidense como puramente binario. La herramienta que con mayor frecuencia evita el cierre administrativo, incluso en un sistema que legalmente lo permite, es la resolución continua (continuing resolution, CR). Una ley de financiación provisional y de corta duración que el Congreso aprueba para mantener el gasto a los niveles del ejercicio anterior mientras continúan las negociaciones sobre el presupuesto definitivo. Una figura que cumple una función muy próxima a la de nuestra propia prórroga, con la diferencia decisiva de que exige un acto legislativo expreso y repetido en lugar de operar de forma automática. Ni siquiera el sistema americano funciona según su lógica más extrema: los cierres administrativos son la excepción cuando la negociación de la CR fracasa.
    
    Existe, además, una tercera vía. El sistema británico de Westminster articula la autorización del gasto mediante un procedimiento de Estimaciones (Estimates) y Votos a Cuenta (Votes on Account), que permite a la Administración seguir operando con normalidad incluso antes de que el Parlamento apruebe la totalidad de las Estimaciones del ejercicio, precisamente porque el sistema de aprobación del gasto británico nunca se articuló en torno a una única ley de presupuestos con fecha límite de caducidad tan rígida como la española o la estadounidense.
    
    Los cierres administrativos más señalados de la historia reciente estadounidense, y el debate que generan
    
    Sin entrar en el detalle de las disputas políticas concretas que los motivaron, conviene retener los tres cierres administrativos más prolongados de la historia estadounidense. El de noviembre de 1995 y diciembre de 1995-enero de 1996, bajo la presidencia de Bill Clinton, con un tramo de 21 días consecutivos. El de octubre de 2013, de 16 días. Y, el más largo de todos, entre diciembre de 2018 y enero de 2019, con 35 días consecutivos de cierre parcial.
    
    El debate doctrinal que estos episodios han generado, tanto en la literatura de ciencia política como en la de administración pública estadounidense, se articula en torno a dos posiciones. Quienes defienden el diseño actual sostienen que el cierre administrativo es la manifestación más pura y más seria del poder constitucional de la bolsa del Congreso. Si el Legislativo pudiera ser ignorado sin coste real por el Ejecutivo en ausencia de apropiación aprobada, el propio control parlamentario sobre el gasto quedaría vaciado de su función disuasoria más elemental. Quienes lo critican, la mayoría en la literatura reciente, señalan que el coste económico y social del cierre recae, de forma desproporcionada, sobre empleados públicos sin ninguna responsabilidad y sobre los ciudadanos que dependen de servicios federales suspendidos, mientras que los verdaderos protagonistas de la negociación no sufren ningún coste comparable; y que, precisamente por eso, el cierre administrativo ha derivado hacia una lógica de toma de rehenes políticos más que hacia una genuina discrepancia sobre cuánto debe gastar el Estado y en qué.
    
    La comparación, en definitiva, es una ventana directa a cómo cada tradición constitucional entiende la propia naturaleza del conflicto entre Ejecutivo y Legislativo sobre el gasto público. La prórroga española asume que ese conflicto, si es grave, debe resolverse por la vía política (censura, disolución, elecciones), dejando a la Administración funcionar mientras tanto. El shutdown estadounidense asume que el propio funcionamiento de la Administración es la pieza que se pone en juego en la negociación, precisamente porque el sistema no ofrece ninguna otra palanca constitucional equivalente para resolver el enfrentamiento entre dos poderes que, a diferencia del caso español, no dependen el uno del otro para seguir existiendo.
\end{specialblock}

Por último, antes del 15 de abril del año siguiente, la AIReF debe emitir el Informe sobre los Presupuestos Iniciales de las AAPP, analizando el grado de cumplimiento de las reglas fiscales de los presupuestos finalmente aprobados para ese año (mientras que el ya mencionado Informe sobre los Proyectos y Líneas Fundamentales de Presupuestos de las AAPP de octubre del año anterior analizaba ex-ante el cumplimiento de las reglas fiscales de los proyectos de presupuestos).

\subsection{Ejecución, con las modificaciones presupuestarias integradas}
En la fase de ejecución el sector público estatal implementa las actuaciones cubiertas por los créditos de los PGE aprobados. Por tanto, es esencial entender la vinculación de los créditos recogidos: ¿hasta qué punto la unidad gestora de fondos públicos debe ceñirse a la finalidad para la que los créditos de los PGE fueron aprobados?

La vinculación se establece a nivel orgánico, funcional y económico. La vinculación orgánica se produce a nivel de servicio o de Organismo Autónomo y supone que el gestor no puede, en principio, trasladar fondos entre servicios o entre Organismos Autónomos. La vinculación funcional, e nivel de programas, no se pueden trasladar fondos de un programa a otro. Por último, la vinculación a nivel económico resulta más compleja, puesto que depende del tipo de gasto y del tipo de centro gestor. Con carácter general, la vinculación es a nivel de concepto.

Sin embargo, la realidad de la ejecución requiere que existan márgenes de flexibilidad. Esta flexibilidad debe ser suficiente para permitir una gestión eficaz., pero no tanto como para convertir en irrelevantes los propios presupuestos aprobados y burlar la voluntad del legislador. La incertidumbre es incluso mayor en el caso de los ingresos, aunque en este caso la propia normativa claro su carácter estimativo. Un presupuesto aprobado en diciembre para regir todo el año siguiente se enfrenta a una realidad que ningún ejercicio de previsión puede anticipar en su totalidad: catástrofes naturales, crisis económicas sobrevenidas, sentencias judiciales inesperadas, cambios normativos con efecto presupuestario inmediato. 

El primer instrumento de flexibilidad es el propio Fondo de Contingencia, pensado para gastos no contemplados que surjan a lo largo del ejercicio, del que puede disponer el Gobierno. El Gobierno también dispone de las modificaciones presupuestarias. que veremos al final del tema. Por último. el Gobierno también puede acotar la ejecución presupuestaria (p.ej. para evitar desviaciones presupuestarias y a avanzado el ejercicio) con un Acuerdo del Cons��io de Ministros de no disponibilidad de crédito o adelantando la Orden de cieITe del Ministro de Hacienda.

Por último, cada gasto se somete a cuatro fases: autori:::ación (acto interno que pone en marcha el procedimiento), compromiso (en el ejemplo de un contrato público, se correspondería con la asignación del contrato al adjudicatario), obligación reconocida ( en el ejemplo citado. se correspondería con l a finalización de la obra o servicio del contrato) y pago (se corresponde con la transferencia de fondos del sector público al adjudicatario).

La rigidez que protege el control parlamentario del gasto entra, en este punto, en tensión directa con la necesidad de gestión eficaz que cualquier Administración real exige. Las modificaciones presupuestarias son, precisamente, la válvula jurídica que intenta resolver esa tensión sin destruir ninguno de los dos valores en juego por completo.

### Las fases del gasto: por qué cuatro pasos y no uno solo

[Seguro] Cada gasto público se somete a cuatro fases sucesivas (ADOP): (i) autorización, acto interno que pone en marcha el procedimiento de gasto, reservando crédito; (ii) disposición o compromiso, la asignación jurídica del gasto a un beneficiario concreto; (iii) reconocimiento de la obligación, constatación de que la contraprestación se ha realizado efectivamente; y, (iv) pago, la salida material de fondos. Esta fragmentación responde a que cada fase constituye un punto de control independiente, sometido a la función interventora de la IGAE, de modo que un error o una irregularidad puede detenerse en cualquiera de los cuatro momentos, antes de que el dinero llegue a salir efectivamente de las arcas públicas. Es la misma lógica de control por capas sucesivas que vemos operar en otros contextos: cuantos más puntos de verificación independientes atraviese una operación antes de consumarse, menor es el riesgo de que un fallo aislado se traduzca en una pérdida efectiva de fondos públicos.

### Las modificaciones presupuestarias: de la más flexible a la más rígida, y por qué en ese orden concreto

Recupero aquí las seis figuras ya trabajadas, pero ordenadas ahora, deliberadamente, de menor a mayor rigidez de aprobación.

\paragraph{Transferencia de crédito}
La Transferencia de crédito es la figura más flexible. Consiste en una autorización ministerial y presidentes de OOAA, con visto bueno de Hacienda. Se reconoce como flexible puesto que no altera el volumen total del gasto, solo lo redistribuye dentro de límites acotados: nunca entre operaciones financieras y no financieras, ni de capital a corrientes, ni entre secciones distintas. 

Ampliaciones y generaciones de crédito.
Las amplicaciones y generaciones de crédito son una autorización del Ministro de Hacienda, sin necesidad de las Cortes, que responden a mecanismos ya previstos de antemano. Por ejemplo, un crédito reconocido como ampliable por ley, o un ingreso inesperado directamente vinculado al gasto que genera.

\paragraph{Incorporación de remanentes} 
La incporación de remanentes es una figura cuya competencia exclusiva es del Ministro de Hacienda. Se puede ejecutar sobre créditos no gastados del ejercicio anterior, ya identificados como incorporables en los propios PGE. 

\paragraph{Créditos extraordinarios y suplementos de crédito}
Por último, los créditos extraordinarios y/o supleentos de crédito son la figura más rígida de todas, reservada a necesidades genuinamente extraordinarias y no previstas. 

Su aprobación exige la participación del Consejo de Ministros, si se financia con cargo al Fondo de Contingencia, o por las Cortes mediante norma con rango de ley en cualquier otro caso.

### El caso que mejor ilustra todo este mecanismo en funcionamiento real: la DANA de Valencia

Pocas veces un episodio ilustra con tanta nitidez el funcionamiento efectivo de estas figuras como la respuesta presupuestaria a la DANA que azotó la Comunitat Valenciana. El Gobierno articuló su respuesta financiera mediante una cascada sucesiva de créditos extraordinarios y suplementos de crédito, todos ellos habilitados por una serie encadenada de Reales Decretos-leyes.\footnote{%
    El RDL 6/2024, de 5 de noviembre, que estableció el marco jurídico habilitante; el RDL 7/2024, de 11 de noviembre, que impulsó el Plan de respuesta inmediata, reconstrucción y relanzamiento, con una disposición adicional específica autorizando al Consejo de Ministros a facilitar la financiación necesaria; y el RDL 1/2025, de 28 de enero, que amplió los mecanismos de financiación disponibles.
} Sobre esta base jurídica, el Consejo de Ministros fue aprobando sucesivos créditos extraordinarios y suplementos de crédito financiados con cargo al Fondo de Contingencia.\footnote{%
    $50$ millones de euros para el Ministerio de Derechos Sociales en diciembre de 2024; $560$ millones de suplemento de crédito para compensar gastos extraordinarios ya ejecutados por la propia Generalitat Valenciana en noviembre de 2025; incluso una partida de $135.000$ euros para reponer vehículos de la Jefatura Central de Tráfico dañados por el temporal, tramitada como crédito extraordinario propio, con cargo al programa presupuestario específicamente creado para la ocasión: Contingencias asociadas a la Depresión Aislada en Niveles Altos (DANA) de 202; un programa nuevo, abierto en cada sección presupuestaria afectada, precisamente para poder rastrear con precisión, dentro de la contabilidad pública, todo el gasto vinculado a esta única catástrofe.
}

En paralelo, la propia Generalitat Valenciana recurrió a su propio mecanismo autonómico de crédito extraordinario, aprobando $700$ millones de euros en diciembre de 2024, y otro de más de $2.364$ millones de euros en marzo de 2025, financiados en parte con cargo al Fondo de Liquidez Autonómico, instrumento de rescate condicionado que se asemeja a una restricción blanda: una Comunidad Autónoma en dificultades recurriendo al mecanismo de auxilio financiero que la teoría de RODDEN identifica como generador de un riesgo de indisciplina si no va acompañado de condicionalidad. En este caso concreto, la condicionalidad se justifica por la naturaleza genuinamente excepcional e imprevisible del gasto, y no por ninguna indisciplina fiscal previa de la Comunidad Autónoma afectada.

### La crítica de fondo: cuando la modificación erosiona el valor informativo del propio presupuesto aprobado

[Suposición] Cierro este epígrafe con una crítica que conviene traer de vuelta desde el Bloque 1, porque encuentra aquí su aplicación más precisa: 

Ya señalamos, al hablar de las debilidades del proceso presupuestario diagnosticadas por la AIReF, que los PGE no se comparan con los gastos e ingresos reales del año anterior, sino con los gastos e ingresos que se presupuestaron inicialmente para dicho año. 

Las modificaciones presupuestarias son el mecanismo técnico que hace posible esa divergencia entre presupuesto inicial y presupuesto finalmente ejecutado. Un Gobierno puede presentar, y las Cortes aprobar, un Proyecto de Presupuestos que cumple con holgura los objetivos de estabilidad presupuestaria en el momento de su aprobación, para después recurrir a créditos extraordinarios, suplementos y ampliaciones que alteran sustancialmente el volumen y la composición final del gasto, con un grado de escrutinio parlamentario sensiblemente inferior al que recibió el propio Proyecto de Ley de Presupuestos, especialmente cuando esos créditos se financian con cargo al Fondo de Contingencia, aprobados directamente por Consejo de Ministros sin pasar en absoluto por las Cortes.

No se trata de sugerir que el recurso a estas figuras sea ilegítimo cuando responde a necesidades extraordinarias, sería difícil imaginar una respuesta más rápida y eficaz, sino de señalar que la misma flexibilidad que permite a un Estado responder con agilidad ante lo verdaderamente imprevisto es la misma flexibilidad que puede terminar vaciando el valor de la aprobación parlamentaria de los PGE.

\subsection{Control: las cuatro perspectivas y los tres poderes}

### Por qué el control no es una única función, sino cuatro entrelazadas

[Probable] Cierro el ciclo presupuestario con la fase que da sentido retroactivo a todas las anteriores. De nada serviría un marco jurídico cuidadosamente diseñado, unos principios bien justificados, y un procedimiento de aprobación cargado de garantías, si nadie verificara que lo aprobado se ha ejecutado tal y como se decidió. 

El control se ejerce desde cuatro perspectivas distintas: legalidad, eficacia/eficiencia, cumplimiento de reglas fiscales, y control político; y por los tres poderes del Estado. Sin embargo, conviene no tratarlas como cuatro compartimentos, pues en la práctica se entrelazan y a veces entran en tensión unas con otras.

\subsubsection{El control de legalidad: dos capas, con un grado de independencia muy distinto entre ellas}
El control interno lo ejerce la IGAE. La IGAE ejerce una función interventora de control previo a la aprobación de cualquier acto de gasto, un control financiero permanente, verificando continuamente el cumplimiento de la normativa contable y los principios de buena gestión financiera, y, por último, de auditoría pública. El control externo corresponde al Tribunal de Cuentas y, en última instancia, al poder judicial. 

Conviene señalar una distinción. La IGAE no es un órgano independiente en el sentido que sí lo es la AIReF. La propia IGAE forma parte de la estructura jerárquica del Ministerio de Hacienda, el mismo departamento cuyo gasto está llamada a fiscalizar. Esta tensión de diseño, presente en cualquier sistema comparado de control interno de la Administración, genere el problema de que el control ejercido carece, por su propia posición orgánica, de la distancia institucional que sí caracteriza a un órgano genuinamente externo como el Tribunal de Cuentas o, en su terreno específico, la AIReF.

Sobre el Tribunal de Cuentas, conviene añadir una crítica de actualidad que se ha puesto de manifiesto de forma reiterada: la lentitud de sus procedimientos jurisdiccionales de responsabilidad contable. En muchas ocasiones documentadas se han prolongado durante más de una década desde los hechos enjuiciados hasta la resolución firme, un problema que plantea la misma pregunta que ya formulamos al hablar de la credibilidad de las sanciones fiscales europeas: ¿qué valor disuasorio real tiene un mecanismo de control cuya respuesta llega muchos años después?

\subsubsection{Control de eficacia y eficiencia: la pieza más joven, y todavía la más incompleta}
Ya trabajamos, con WILDAVSKY, por qué este tipo de control resulta estructuralmente más difícil que el control de legalidad. 

IGAE y Tribunal de Cuentas pueden ejercerlo, pero suelen concentrar su actividad en el control de legalidad. 

La comparación internacional resulta aquí especialmente instructiva. Instituciones fiscalizadoras superiores de otros países llevan décadas desarrollando auditorías específicas de value for money, evaluando sistemáticamente si un programa público concreto logra sus objetivos declarados al menor coste posible, siguiendo los estándares internacionales de auditoría del desempeño que promueve INTOSAI (la organización internacional de instituciones fiscalizadoras superiores). España, con los \textit{spending reviews} de la AIReF, se ha incorporado tarde y de forma todavía parcial a esta tradición. Limitada, además, a los encargos puntuales que las distintas Administraciones deciden solicitarle, y no a una función de auditoría del desempeño sistemática y universal como la que sí ejercen de oficio otras instituciones fiscalizadoras comparadas.

\subsubsection{Control del cumplimiento de las reglas fiscales: cuando el árbitro técnico y el jugador discrepan en público}

Ya trabajamos la arquitectura formal: informes de Hacienda/IGAE dos veces al año, los dos grandes informes anuales de la AIReF, notificación semestral a Eurostat, evaluación de la Comisión Europea en mayo y noviembre. 

Conviene subrayar algo que la descripción puramente procedimental esconde. Esta relación de control no siempre transcurre con la fluidez que el diseño institucional presupone. 

La AIReF se define a sí misma como una institución que actúa con total y plena independencia orgánica y funcional respecto de cualquier entidad pública o privada, adscrita al Ministerio de Hacienda únicamente a efectos organizativos y presupuestarios. Pero esa independencia formal no ha impedido episodios de fricción real y públicamente documentados. La propia AIReF ha denunciado no haber recibido a tiempo la información necesaria para evaluar correctamente el Programa de Estabilidad del Gobierno, llegando incluso a posponer la publicación de su propio informe preceptivo por falta de datos, y ha discrepado abiertamente de las previsiones macroeconómicas del propio Ejecutivo, señalando un optimismo gubernamental sistemático sobre el crecimiento del PIB y sobre la recaudación esperada. Más allá de la fricción técnica puntual, ha llegado a plantearse una discusión sobre la propia independencia institucional de la AIReF frente a eventuales intentos de influencia política en el nombramiento de sus responsables. 

Un debate que confirma que la tensión entre credibilidad técnica y control político de las AAIs en general.

\subsubsection{Control político: el más antiguo de los cuatro, y el que cierra el círculo con el que empezamos}

Ejercido por las Cortes —mociones, interpelaciones, comisiones de investigación, solicitudes de comparecencia— y, en un sentido más amplio, por la opinión pública a través de medios de comunicación y servicios de estudios.

Este es el control más antiguo de los cuatro, el mismo control que las Cortes de Cádiz ejercían sobre la Contaduría Mayor de Cuentas, formulado ahora con instrumentos parlamentarios modernos. Sin embargo, presenta una limitación estructural evidente. Cuando el Gobierno cuenta con mayoría parlamentaria sólida, los mecanismos de control político rara vez producen consecuencias reales, porque es la propia mayoría que sostiene al Gobierno la que decide si esos mecanismos prosperan o se archivan. 

Pero precisamente en los periodos de mayor fragmentación parlamentaria es cuando el control político parlamentario recupera, paradójicamente, un vigor mayor. Un Gobierno en minoría no puede blindarse frente al escrutinio de una oposición con capacidad real de bloquear, enmendar o cuestionar públicamente su gestión presupuestaria. La propia comparecencia reciente de responsables de la AIReF ante comisiones parlamentarias donde la oposición cuenta con mayoría es un ejemplo. El control técnico y el control político se refuerzan mutuamente cuando el Parlamento es quien convoca y escucha a la institución técnica independiente.

## 2.5. Cierre: las clasificaciones, solo como anuncio

Al hilo de cómo se elabora, se ejecuta y se controla el presupuesto, conviene cerrar con una última precisión que no desarrollamos aquí en detalle: los ingresos y los gastos que hemos visto recorrer todo este ciclo se ordenan, a efectos de su elaboración y de su análisis, según tres criterios distintos —orgánico (quién gasta o ingresa), económico (cómo se gasta o se ingresa) y funcional o por programas (en qué se gasta) —. No los desarrollamos aquí con las cifras y el detalle que merecen: ese es, precisamente, el objeto propio del Tema 22, al que remitimos expresamente para su aplicación completa sobre las grandes cifras de los PGE.






















































































\clearpage
\section{Conclusión} 


\subsection{Recapitulación}


\subsection{Extensiones y relación con otras partes del Temario}



\subsection{Opinión}

\subsubsection{Idea final (Salida o cierra)}

\clearpage
\pagenumbering{gobble}
\MyBibliography



\MyBibItem{DAWID y GATTI}{Chapter 2 - Agent-Based Macroeconomics}{2018}{https://doi.org/10.1016/bs.hescom.2018.02.006}

- Constitución Política de la Monarquía Española (1812)
- Decreto de las Cortes de 22 de marzo de 1811 (inicio de la formación de presupuestos).
- Decreto de 7 de agosto de 1813 (Reglamento de la Contaduría Mayor de Cuentas).
- Estatuto Real de 1834.
- Decreto de 1849 (instauración de la contabilidad por ejercicios anuales) y Ley de Contabilidad y Administración Pública de 1850.
- Ley de Administración y Contabilidad de la Hacienda Pública, de 1 de julio de 1911.
- Ley 11/1977, de 4 de enero, General Presupuestaria.
- Real Decreto Legislativo 1091/1988, de 23 de septiembre, Texto Refundido de la Ley General Presupuestaria.
- Ley 47/2003, de 26 de noviembre, General Presupuestaria (norma vigente).
- ARTOLA, M. (1986): La Hacienda del siglo XIX, referencia sobre la reforma contable de 1849-1850.
- RODRÍGUEZ BEREIJO, A.: doctrina sobre la institución presupuestaria como sometimiento del poder financiero al imperio de la Ley (citada en el Voto particular de la STC 116/1994).
- Constitución Española de 1978, arts. 134, 135 y 136 (texto vigente tras la reforma de 2011).
- Reforma del artículo 135 de la Constitución Española, de 27 de septiembre de 2011 (BOE nº 233, de 27 de septiembre de 2011).
- Reforma del artículo 13.2 de la Constitución Española, de 27 de agosto de 1992 (BOE nº 207, de 28 de agosto de 1992) — primera reforma constitucional española, como contraste procedimental.
- GARCÍA-ESCUDERO MÁRQUEZ, P. (2012): "La acelerada tramitación parlamentaria de la reforma del artículo 135 de la Constitución", Teoría y Realidad Constitucional, núm. 29, UNED.
- GARCÍA-ESCUDERO MÁRQUEZ, P. (2009): "La reforma del artículo 135: ¿son suficientes trece días para la tramitación parlamentaria de una reforma constitucional?", Cuadernos de Derecho Público, núm. 38, INAP.
- BAR CENDÓN, A. (2012): "La reforma constitucional y la gobernanza económica de la Unión Europea", Teoría y Realidad Constitucional, núm. 30, UNED.
- Enmienda a la totalidad de Izquierda Unida-Iniciativa per Catalunya Verds a la Proposición de reforma del art. 135 CE, 1 de septiembre de 2011 (Congreso de los Diputados).
- STC 248/2007, de 13 de diciembre (contenido constitucionalmente admisible de la Ley de Presupuestos).
- Grundgesetz alemana, reforma de 2009 (introducción de la Schuldenbremse o freno constitucional al endeudamiento), como referencia comparada de la reforma española.
- Tratado de Estabilidad, Coordinación y Gobernanza en la Unión Económica y Monetaria (Fiscal Compact), firmado el 2 de marzo de 2012.
- Ley 47/2003, de 26 de noviembre, General Presupuestaria.
- Ley Orgánica 2/2012, de 27 de abril, de Estabilidad Presupuestaria y Sostenibilidad Financiera.
- Ley Orgánica 2/2012, de 27 de abril, de Estabilidad Presupuestaria y Sostenibilidad Financiera, arts. 3 a 13 y 26.
- Tratado de Maastricht (Tratado de la Unión Europea), 1992, Protocolo sobre el procedimiento aplicable en caso de déficit excesivo.
- TIETMEYER, H. (2005): memorias sobre la negociación del Tratado de Maastricht y el origen de los valores de referencia del 60% y el 3%.
- ALESINA, A. y TABELLINI, G.: literatura sobre ciclos políticos presupuestarios y descuento temporal excesivo de gobiernos sujetos a ciclo electoral.
- International Budget Partnership: Open Budget Survey, ediciones sucesivas (comparación internacional de transparencia presupuestaria).
- Fondo Monetario Internacional: Fiscal Transparency Code y guías de buenas prácticas sobre fondos extrapresupuestarios.
- Reforma del marco de gobernanza económica de la Unión Europea, en vigor desde 2024 (Reglamentos (UE) 2024/1263 y 2024/1264), sobre el papel central de la regla de gasto.
Comisión Taft sobre Economía y Eficiencia (1910-1912), EEUU.
CLEVELAND, F. A. (1915): "The Evolution of the Budget Idea in the United States", The Annals of the American Academy of Political and Social Science.
WILLOUGHBY, W. F. (1918): The Movement for Budgetary Reform in the States.
Budget and Accounting Act de 1921 (EEUU).
WILDAVSKY, A. (1964): The Politics of the Budgetary Process, Little, Brown.
KYDLAND, F. E. y PRESCOTT, E. C. (1977): "Rules Rather than Discretion: The Inconsistency of Optimal Plans", Journal of Political Economy, vol. 85, núm. 3.
OATES, W. E. (1972): Fiscal Federalism, Harcourt Brace Jovanovich.
KORNAI, J. (1980): Economics of Shortage (origen del concepto de restricción presupuestaria blanda).
RODDEN, J. (2006): Hamilton's Paradox: The Promise and Peril of Fiscal Federalism, Cambridge University Press.
Bill of Rights inglesa de 1689.
Constitución prusiana de 1850; conflicto constitucional prusiano de 1862-1866; Ley de Indemnidad prusiana de 3 de septiembre de 1866.
JÈZE, G.: Cours de science des finances et de législation financière française (ediciones sucesivas desde comienzos del siglo XX).
LEROY-BEAULIEU, P. (1877): Traité de la science des finances.
- GARCÍA DE ENTERRÍA, E. y FERNÁNDEZ, T.-R.: Curso de Derecho Administrativo, doctrina clásica sobre Administración Institucional y personificación instrumental.
- Ley 40/2015, de 1 de octubre, de Régimen Jurídico del Sector Público, en especial arts. 84 a 127 (Título II, sector público institucional).
- Ley 28/2006, de 18 de julio, de Agencias Estatales para la mejora de los servicios públicos (derogada por la Ley 40/2015).
- Constitución Española de 1978, Título VIII (arts. 137-158), y en especial arts. 142 y 156.
- Ley Orgánica 8/1980, de 22 de septiembre, de Financiación de las Comunidades Autónomas (LOFCA).
- Ley 22/2009, de 18 de diciembre, del sistema de financiación de las CCAA de régimen común.
- Constitución Española de 1978, Disposición Adicional Primera (derechos históricos de los territorios forales).
- Ley 12/2002, de 23 de mayo, por la que se aprueba el Concierto Económico con la Comunidad Autónoma del País Vasco.
- Ley 28/1990, de 26 de diciembre, por la que se aprueba el Convenio Económico entre el Estado y la Comunidad Foral de Navarra.
- Real Decreto Legislativo 2/2004, de 5 de marzo, Texto Refundido de la Ley Reguladora de las Haciendas Locales.
Comunicación de la Comisión Europea, "Reforzar la coordinación de las políticas económicas" (creación del Semestre Europeo), 2010.
Six-Pack: Reglamentos (UE) 1173/2011, 1174/2011, 1175/2011, 1176/2011 y 1177/2011 del Parlamento Europeo y del Consejo, y Directiva 2011/85/UE del Consejo, de 8 de noviembre de 2011, sobre los requisitos aplicables a los marcos presupuestarios de los Estados miembros.
Two-Pack: Reglamentos (UE) 472/2013 y 473/2013 del Parlamento Europeo y del Consejo, de 21 de mayo de 2013 (supervisión ex ante de los proyectos de planes presupuestarios de la eurozona).
Tratado de Estabilidad, Coordinación y Gobernanza en la Unión Económica y Monetaria (Fiscal Compact), 2012.
AIReF: informes sucesivos sobre el Programa de Estabilidad de España, críticos con su falta de desagregación por subsectores y con la ausencia de corresponsabilidad de las Administraciones Territoriales en su elaboración.
Ley Orgánica 2/2012, art. 15, sobre la composición y el reparto de voto del Consejo de Política Fiscal y Financiera.
Constitución de la V República francesa, de 4 de octubre de 1958, art. 40.
CARCASSONNE, G.: doctrina sobre el parlementarisme rationalisé y la caracterización del Parlamento francés "bajo tutela" en materia financiera.
Documents pour servir à l'histoire de l'élaboration de la Constitution du 4 octobre 1958, La Documentation française, 1991 (intervención de Gilbert Devaux ante la comisión constitucional del Consejo de Estado francés sobre el origen del art. 40).
ALESINA, A. y PEROTTI, R. (1996): "Budget Deficits and Budget Institutions", IMF Staff Papers / NBER Working Paper.
ZALMA (1985): "L'hégémonie du ministre des Finances dans le droit budgétaire de l'État", Revue du Droit Public, núm. 6 (crítica comparada al protagonismo del Ministerio de Hacienda en el procedimiento presupuestario francés, de aplicación análoga al caso español).
- Ley 47/2003, General Presupuestaria, arts. 46-58 (modificaciones presupuestarias) y art. 50 (Fondo de Contingencia).
- Real Decreto-ley 6/2024, de 5 de noviembre, de medidas urgentes de respuesta ante los daños causados por la DANA.
- Real Decreto-ley 7/2024, de 11 de noviembre, por el que se adoptan medidas urgentes para el impulso del Plan de respuesta inmediata, reconstrucción y relanzamiento frente a los daños causados por la DANA.
- Real Decreto-ley 1/2025, de 28 de enero, por el que se aprueban medidas urgentes en materia económica, de transporte, de Seguridad Social y para hacer frente a situaciones de vulnerabilidad.
- Decreto-ley 16/2024, de 17 de diciembre, del Consell (crédito extraordinario de 700 millones de euros de la Generalitat Valenciana).
- Decreto-ley 5/2025, de 11 de marzo, del Consell (crédito extraordinario de 2.364,28 millones de euros de la Generalitat Valenciana).
- Notas de prensa del Ministerio de Hacienda, Consejo de Ministros, diciembre de 2024 a abril de 2025, sobre aplicaciones sucesivas del Fondo de Contingencia para gastos derivados de la DANA.
- Ley Orgánica 6/2013, de 14 de noviembre, de creación de la Autoridad Independiente de Responsabilidad Fiscal.
- Directiva 2011/85/UE del Consejo, de 8 de noviembre de 2011 (base de la transposición de la AIReF).
- Informes y comunicados públicos de la AIReF sobre discrepancias con las previsiones macroeconómicas y fiscales del Gobierno en la elaboración del Programa de Estabilidad (episodios recurrentes, incluidos los más recientes).
- National Audit Office (Reino Unido): metodología de auditorías de value for money, como referencia comparada de control de eficiencia del gasto público.
- INTOSAI (Organización Internacional de Entidades Fiscalizadoras Superiores): estándares internacionales de auditoría del desempeño (performance auditing).


% ANEXOS
\clearpage










\end{document}